\chapter{Detección y Evaluación del Modelo (OE2)}
\label{cap:deteccion}

% Ruta de las figuras y de los fragmentos que deja el cuaderno de este capítulo.
\newcommand{\figDET}{capitulos/05_oe2_detector/figures}

% ===================================================================
\section{Introducción}
% ===================================================================

Este capítulo presenta los resultados alcanzados para el segundo objetivo
específico:

\begin{quote}
    \textbf{OE2.} Implementar y evaluar un detector que reciba un espectrograma y
    devuelva cajas tiempo--frecuencia con especie, tipo de llamada y confianza,
    comparándolo con una línea base establecida.
\end{quote}

Los cinco resultados comprometidos en la Tabla~\ref{tab:iov-oe2} ocupan una
sección cada uno: el diseño experimental, la arquitectura propuesta y las líneas
base (RE2.1, sección~\ref{sec:re21}), el \textit{pipeline} reproducible de
entrenamiento, evaluación e inferencia (RE2.2, sección~\ref{sec:re22}), el
informe comparativo de resultados (RE2.3, sección~\ref{sec:re23}), el modelo
final cargable y su demostración (RE2.4, sección~\ref{sec:re24}) y la
publicación del código y del modelo (RE2.5, sección~\ref{sec:re25}). Cada
sección sigue la estructura del Capítulo~\ref{cap:curacion}: el resultado, cómo
se alcanzó, cómo reproducirlo y la verificación de su indicador. La
sección~\ref{sec:discusion-oe2} discute los cinco en conjunto.

El conjunto sobre el que se entrena y se mide es el que produce el
Capítulo~\ref{cap:curacion}: \nWindowsTrain{} ventanas de entrenamiento,
\nWindowsVal{} de validación y \nWindowsTest{} de prueba, repartidas por
grabación y con \nClasses{} clases de especie y tipo de llamada. El costo de
revisión que genera el modelo elegido aquí es materia de OE3.

% ===================================================================
\section{Diseño experimental, arquitectura propuesta y líneas base (RE2.1)}
\label{sec:re21}
\label{sec:arquitectura}
% ===================================================================

\rotulo{Resultado.} La arquitectura propuesta (un AST-Deformable-DETR de
dimensión \mDim{}, \mQueries{} consultas y \mLevels{} niveles, implementada en
\file{src/models/}), sus hiperparámetros, las tres líneas base contra las que
se compara, las recetas de entrenamiento y el protocolo de comparación con su
regla de decisión fijada antes de correr el experimento.

\rotulo{Cómo se alcanzó.} Aplicando la Fase~4 de CRISP-DM
(sección~\ref{sec:fase4}), que fija la forma de la salida, las configuraciones
que se comparan y la regla de decisión antes de entrenar. El punto de partida es
la formulación como detección de objetos sobre el espectrograma
(sección~\ref{sec:objetos}) y la representación que fija el
Capítulo~\ref{cap:curacion}: el modelo recibe exactamente el tensor que
el caché almacena y devuelve cajas en el mismo marco normalizado en el que se
anotaron. Las tres decisiones que lo separan de un detector de imágenes estándar
(extractor preentrenado en audio, pirámide sobre una rejilla de tokens sin
jerarquía y cabeza de predicción de conjuntos) responden cada una a una
propiedad del material documentada en el capítulo anterior.

\rotulo{Cómo reproducirlo.} \file{src/train.py} instancia la propuesta o
cualquiera de las líneas base según la arquitectura que se le indique, con la
receta de entrenamiento que fija esta sección, y copia la configuración completa al
\file{config.json} de la corrida. Así, cada modelo de este capítulo se vuelve a
construir con los mismos hiperparámetros.

\rotulo{Medio de verificación.} Esta sección, que es el documento de diseño, y el
módulo de arquitecturas \file{src/models/} del repositorio del proyecto
(\url{\rRepoUrl}), en el \textit{commit} \rCommit{}; la
sección~\ref{sec:verificacion-re21} contrasta cada elemento del indicador con el
código. La configuración completa de las cuatro corridas comparadas está en la
carpeta RE2.1 de los medios de verificación (Anexo~\ref{anx:medios}).

\subsection{La tarea}

La tarea es la del nivel de ventana de la sección~\ref{sec:formulacion}. La
entrada es una ventana de \clipLen{}\,s convertida en espectrograma mel de
potencia $X \in \mathbb{R}^{1 \times \nMels \times \nFrames}$
(sección~\ref{sec:representacion}), y la salida, el conjunto de la
ecuación~\eqref{eq:modelo}: cajas en el marco normalizado de la ventana, cada una
con una de las \nClasses{} clases y una confianza.

La formulación separa tres preguntas, y el protocolo de la
sección~\ref{sec:protocolo-oe2} las puntúa por separado: si la llamada se
\emph{encontró} (detección), si la caja quedó bien \emph{encuadrada}
(localización) y si la etiqueta es \emph{correcta} (clasificación). Mantenerlas
separadas no es un refinamiento metodológico: es lo que permite decir, al final
del capítulo, dónde está el error.

\subsection{Arquitectura propuesta: AST-Deformable-DETR}

La Figura~\ref{fig:arquitectura} recorre el modelo de la entrada a la salida.

\begin{figure}[htbp]
    \centering
    \caption{Arquitectura propuesta}
    \label{fig:arquitectura}
    \resizebox{\textwidth}{!}{%
        \begin{tikzpicture}[node distance=4mm and 6mm, every node/.style={font=\footnotesize}]
            \node[dato] (mel) {mel\\$1 \times \nMels \times \nFrames$};
            \node[etapa, right=of mel] (fe) {compresión\\log-mel};
            \node[etapa, right=of fe] (bn) {BatchNorm\\$\cdot\,1/2$};
            \node[etapa, right=of bn, fill=celeste!15, draw=celeste!60!black] (ast)
            {AST\\parches $16 \times 16$\\paso $10 \times \mTimeStride$\\\mAstLayers{} capas, \mHidden-d};
            \node[dato, right=of ast] (tok)
            {tokens\\$\mTokensFreq \times \mTokensTime \times \mHidden$};
            \node[etapa, below=14mm of tok] (proj) {proyección\\$\mHidden \to \mDim$};
            \node[etapa, left=of proj] (pyr) {pirámide\\$4\times, 2\times, 1\times, \tfrac{1}{2}\times$};
            \node[etapa, left=of pyr, fill=celeste!15, draw=celeste!60!black] (dec)
            {decodificador deformable\\\mDecoderLayers{} capas, \mQueries{} consultas\\\mHeads{} cabezas, \mPoints{} puntos};
            \node[dato, left=of dec] (out) {\mQueries{} $\times$ (\nClasses{} logits\\+ caja $c_x c_y w h$)};
            \draw[flecha] (mel) -- (fe);
            \draw[flecha] (fe) -- (bn);
            \draw[flecha] (bn) -- (ast);
            \draw[flecha] (ast) -- (tok);
            \draw[flecha] (tok) -- (proj);
            \draw[flecha] (proj) -- (pyr);
            \draw[flecha] (pyr) -- (dec);
            \draw[flecha] (dec) -- (out);
            \node[nota, above=1mm of ast] {preentrenado en AudioSet, se afina entero};
            \node[nota, below=1mm of dec] {desde cero, o preentrenado con aves};
        \end{tikzpicture}}
\end{figure}

El mel entra por la izquierda, atraviesa el extractor y vuelve, en la fila
inferior, como \mQueries{} propuestas de caja con sus logits de clase. Tres
decisiones definen ese recorrido.

\rotulo{Un extractor preentrenado en audio, no en imágenes.} El AST
(\emph{Audio Spectrogram Transformer}) preentrenado en AudioSet
(\mAstCheckpoint{}) sustituye al \textit{backbone} convolucional de imágenes que
traen los detectores estándar. La brecha de modalidad entre imagen y audio
justifica esa sustitución \citep{zhu_sound_2025}: las estadísticas de un
espectrograma (ejes asimétricos, sin invariancia a la traslación en
frecuencia) no son las de una fotografía.

\rotulo{Una pirámide multiescala sobre los tokens.} El AST produce una única
rejilla de $\mTokensFreq \times \mTokensTime$ tokens, es decir, una sola escala,
mientras que la duración mediana de las clases va de \mShortestClassMs{}\,ms
(\mShortestClass{}) a \mLongestClassS{}\,s (\mLongestClass{}), un rango de más de
dos órdenes de magnitud (sección~\ref{sec:geometria}). La pirámide corrige que
las arquitecturas de tipo \textit{Vision Transformer} (ViT) produzcan representaciones de escala única
\citep{zhu_sound_2025}, siguiendo la receta de ViTDet: deconvoluciones y una
convolución con paso 2 sobre el mismo mapa.

\rotulo{Un decodificador deformable de predicción de conjuntos.} La cabeza es el
decodificador de Deformable DETR \citep{zhu_deformable_2021}: \mQueries{}
consultas aprendidas atienden a unos pocos puntos muestreados alrededor de una
caja de referencia, en los \mLevels{} niveles, y refinan la caja capa a capa. No
hay anclas, y el emparejamiento húngaro durante el entrenamiento fuerza una
predicción por evento.

\subsubsection{Compresión y normalización de entrada}

El caché guarda el mel de potencia sin comprimir, y la primera capa del modelo lo
comprime con $\log(\mathrm{mel} + 10^{-6})$: el log-mel, la representación sobre
la que se preentrenó el AST (sección~\ref{sec:frontends}).
A continuación, una \code{BatchNorm2d} sin parámetros afines lleva la entrada a
media 0 y varianza 1 por lote, y se divide entre 2: es el rango con el que el AST
se preentrenó. La corrida que se compara en la sección~\ref{sec:re23} usa
\mFrontend{}.

\subsubsection{Extractor: AST con paso temporal adaptado}

El AST corta el espectrograma en parches de $16 \times 16$ y los proyecta a
\mHidden{} dimensiones. Con el paso original de 10 en ambos ejes, una ventana de
\nMels{} bandas por \nFrames{} tramas produce
$\mTokensFreq \times \mTokensTime$ tokens, es decir unos \mTokenMs{}\,ms de audio
por token en el eje temporal. La incrustación posicional del \emph{checkpoint},
entrenada para otra rejilla, se interpola bilinealmente a la nueva. El extractor
se afina entero, a la misma tasa que la cabeza: con el extractor congelado la
propuesta quedó última en una comparación anterior, y descongelarlo es lo que la
pone a la altura de las líneas base.

\subsubsection{Pirámide multiescala}

Los tokens se proyectan de \mHidden{} a \mDim{} dimensiones y se reordenan como
un mapa $\mTokensFreq \times \mTokensTime$. Sobre ese mapa, cuatro ramas producen
los \mLevels{} niveles de la Figura~\ref{fig:piramide}.

\begin{figure}[htbp]
    \centering
    \caption{Pirámide multiescala}
    \label{fig:piramide}
    \begin{tikzpicture}[x=0.40mm, y=0.40mm, every node/.style={font=\footnotesize}]
        \draw[fill=black!8, draw=black!60] (150, 95) rectangle (182, 107);
        \node[anchor=south] at (166, 108)
        {tokens $\mTokensFreq \times \mTokensTime$, \mDim{} canales};
        \draw[fill=azul!12, draw=azul!70] (0, 0) rectangle (128, 48);
        \draw[fill=azul!12, draw=azul!70] (150, 0) rectangle (214, 24);
        \draw[fill=azul!12, draw=azul!70] (236, 0) rectangle (268, 12);
        \draw[fill=azul!12, draw=azul!70] (290, 0) rectangle (306, 6);
        \node[anchor=north, text=azul!80!black, align=center] at (64, -3)
        {$4\times$: $48 \times 128$\\\mCellZero{}\,ms por celda};
        \node[anchor=north, text=azul!80!black, align=center] at (182, -3)
        {$2\times$: $24 \times 64$\\\mCellOne{}\,ms};
        \node[anchor=north, text=azul!80!black, align=center] at (252, -3)
        {$1\times$: $12 \times 32$\\\mCellTwo{}\,ms};
        \node[anchor=north, text=azul!80!black, align=center] at (298, -3)
        {$\tfrac{1}{2}\times$: $6 \times 16$\\\mCellThree{}\,ms};
        \draw[flecha] (152, 95) -- (64, 50) node[pos=0.5, above left, nota] {2 deconv $2\times$};
        \draw[flecha] (162, 95) -- (182, 26) node[pos=0.6, left, nota] {deconv $2\times$};
        \draw[flecha] (172, 95) -- (252, 14) node[pos=0.5, above right, nota] {identidad};
        \draw[flecha] (180, 95) -- (298, 8) node[pos=0.55, above right, nota] {conv paso 2};
    \end{tikzpicture}
\end{figure}

Cada rama cambia la resolución del mismo mapa: dos deconvoluciones la
cuadruplican, una la duplica, la identidad la conserva y una convolución con
paso~2 la reduce a la mitad. La Tabla~\ref{tab:niveles} recoge el tamaño de celda
de cada nivel.

\begin{table}[htbp]
    \centering
    \caption{Niveles de la pirámide}
    \label{tab:niveles}
    \small
    \begin{tabular}{@{}lccc@{}}
        \toprule
        \textbf{Escala} & \thead{Mapa\\(bandas $\times$ tramas)} & \thead{ms por celda} & \thead{Bandas mel\\por celda} \\
        \midrule
        \input{\figDET/tablas/niveles} \\
        \bottomrule
    \end{tabular}
\end{table}

El nivel más fino resuelve \mCellZero{}\,ms por celda, por debajo de la mediana
de la clase más corta (\mShortestClass{}, \mShortestClassMs{}\,ms); el más grueso,
\mCellThree{}\,ms, es del orden de las sílabas largas. La
Figura~\ref{fig:duraciones} sitúa las \nClasses{} clases frente a esas cuatro
resoluciones y frente a la ventana.

\begin{figure}[htbp]
    \centering
    \caption{Duración de las clases frente a la pirámide}
    \label{fig:duraciones}
    \begin{tikzpicture}
        \begin{axis}[
                width=0.92\textwidth, height=11cm,
                xmode=log, xmin=40, xmax=30000,
                xlabel={duración (ms, escala logarítmica)},
                ytick=data, yticklabels from table={\figDET/data/duraciones.dat}{clase},
                yticklabel style={font=\scriptsize\ttfamily},
                ymin=-1, ymax=25,
                axis line style={black!50}, tick style={black!50},
                xmajorgrids, grid style={black!12},
                log ticks with fixed point, clip=false,
                x tick label style={/pgf/number format/1000 sep={\,}},
            ]
            \addplot[only marks, mark=*, mark size=1.6pt, azul,
                error bars/.cd, x dir=both, x explicit, error bar style={azul!70}]
            table[x=med, y expr=\coordindex,
                    x error minus expr=\thisrow{med}-\thisrow{q1},
                    x error plus expr=\thisrow{q3}-\thisrow{med}]
                {\figDET/data/duraciones.dat};
            \draw[celeste!80!black, dashed] (axis cs:\mCellZeroPlain,-1) -- (axis cs:\mCellZeroPlain,25)
            node[pos=1, above, rotate=90, anchor=west, font=\scriptsize] {$4\times$ (\mCellZero{}\,ms)};
            \draw[celeste!80!black, dashed] (axis cs:\mCellOnePlain,-1) -- (axis cs:\mCellOnePlain,25)
            node[pos=1, above, rotate=90, anchor=west, font=\scriptsize] {$2\times$};
            \draw[celeste!80!black, dashed] (axis cs:\mCellTwoPlain,-1) -- (axis cs:\mCellTwoPlain,25)
            node[pos=1, above, rotate=90, anchor=west, font=\scriptsize] {$1\times$};
            \draw[celeste!80!black, dashed] (axis cs:\mCellThreePlain,-1) -- (axis cs:\mCellThreePlain,25)
            node[pos=1, above, rotate=90, anchor=west, font=\scriptsize] {$\tfrac{1}{2}\times$ (\mCellThree{}\,ms)};
            \draw[rosa!80!black, dashed] (axis cs:3000,-1) -- (axis cs:3000,25)
            node[pos=1, above, rotate=90, anchor=west, font=\scriptsize] {ventana (\clipLen{}\,s)};
        \end{axis}
    \end{tikzpicture}
\end{figure}

Las líneas punteadas son las cuatro resoluciones y la ventana. Las duraciones
medianas de casi todas las clases caen entre la resolución más fina y la ventana,
y solo \mClassesOverWindow{} clases (las dos de ventana) la superan.

\subsubsection{Decodificador deformable}

El decodificador tiene \mDecoderLayers{} capas idénticas
(Figura~\ref{fig:decoder}), de dimensión \mDim{}, con \mHeads{} cabezas de
atención, FFN de \mFfn{} y \textit{dropout} \mDropout{}. Cada una de las
\mQueries{} consultas lleva dos vectores aprendidos, contenido y posición, y de
la posición se deriva un punto de referencia inicial $(c_x, c_y)$ que, con un
tamaño inicial de 0,1, forma la caja de referencia de la primera capa. En cada
capa ocurre lo siguiente:

\begin{enumerate}[nosep]
    \item \textbf{Autoatención} entre consultas, con la posición sumada a las
          claves y consultas en cada capa, para que las consultas no se diluyan y
          puedan repartirse los eventos entre sí.
    \item \textbf{Atención cruzada deformable} (Figura~\ref{fig:deformable}): en
          lugar de atender a todas las celdas de los cuatro mapas, cada consulta
          y cada cabeza muestrean \mPoints{} puntos por nivel, en posiciones que
          la propia consulta predice como desplazamientos relativos al tamaño de
          su caja de referencia, y los combinan con pesos que también predice
          (\textit{softmax} sobre los $\mLevels \times \mPoints$ puntos). Los
          desplazamientos arrancan en una rejilla radial y no en ruido, como en
          el diseño original \citep{zhu_deformable_2021}. Es lo que hace tratable
          atender a un mapa de $48 \times 128$ celdas y lo que permite mezclar
          escalas sin elegir una.
    \item \textbf{FFN} y normalización residual.
    \item \textbf{Cabezas por capa.} Una cabeza lineal produce los \nClasses{}
          logits de clase y un perceptrón multicapa (MLP) de tres capas produce un delta de caja que se
          suma, en el espacio logit, a la caja de referencia; la caja refinada es
          la referencia de la capa siguiente, sin gradiente entre capas. Cada capa
          se supervisa por separado (pérdidas auxiliares).
\end{enumerate}

\begin{figure}[htbp]
    \centering
    \begin{minipage}[t]{0.47\textwidth}
        \centering
        \captionof{figure}{Una capa del decodificador}
        \label{fig:decoder}
        \resizebox{\textwidth}{!}{%
            \begin{tikzpicture}[node distance=2.4mm and 6mm, every node/.style={font=\footnotesize}]
                \node[dato] (q) {consultas $q$ ($\mQueries \times \mDim$)};
                \node[etapa, below=of q] (sa) {autoatención sobre $q + p$};
                \node[etapa, below=of sa] (n1) {suma y norma};
                \node[etapa, below=of n1, fill=celeste!15, draw=celeste!60!black] (ca)
                {atención cruzada deformable};
                \node[etapa, below=of ca] (n2) {suma y norma};
                \node[etapa, below=of n2] (ffn) {FFN $\mDim \to \mFfn \to \mDim$};
                \node[etapa, below=of ffn] (n3) {suma y norma};
                \node[etapa, below=of n3, xshift=-20mm] (cls) {cabeza de clase\\\nClasses{} logits};
                \node[etapa, below=of n3, xshift=20mm] (box) {cabeza de caja\\$\Delta b$};
                \node[dato, below=of box] (ref) {$b \leftarrow \sigma(\sigma^{-1}(b) + \Delta b)$};
                \node[dato, right=of ca] (maps) {mapas de la pirámide\\\mLevels{} niveles};
                \node[dato, left=of ca] (pos) {posición $p$\\y caja de referencia $b$};
                \foreach \a/\b in {q/sa, sa/n1, n1/ca, ca/n2, n2/ffn, ffn/n3, box/ref}
                \draw[flecha] (\a) -- (\b);
                \draw[flecha] (n3) -- (cls);
                \draw[flecha] (n3) -- (box);
                \draw[flecha] (maps) -- (ca);
                \draw[flecha] (pos) -- (ca);
                \draw[flecha] (pos.north) |- (sa.west);
                \draw[flecha, dashed] (ref.east) -- ++(9mm, 0) |- (q.east)
                node[pos=0.25, right, nota, align=left] {referencia de la\\capa siguiente};
                \node[nota, below=1mm of ref] {$\times$ \mDecoderLayers{} capas, cada una supervisada};
            \end{tikzpicture}}
    \end{minipage}\hfill
    \begin{minipage}[t]{0.49\textwidth}
        \centering
        \captionof{figure}{Muestreo deformable de una cabeza sobre un nivel}
        \label{fig:deformable}
        \resizebox{\textwidth}{!}{%
            \begin{tikzpicture}[x=7mm, y=7mm, every node/.style={font=\footnotesize}]
                \draw[step=7mm, black!15] (0, 0) grid (12, 6);
                \draw[black!50] (0, 0) rectangle (12, 6);
                \node[anchor=south west, nota] at (0, 6.1)
                {un nivel de la pirámide (celdas de \mDim{} canales)};
                \draw[rosa!80!black, thick] (3.5, 1.5) rectangle (8.5, 4.5);
                \fill[rosa!80!black] (6, 3) circle (2pt)
                node[below, text=rosa!80!black, yshift=-1pt] {$(c_x, c_y)$};
                \draw[<->, rosa!80!black] (3.5, 0.9) -- (8.5, 0.9) node[midway, below] {$w$};
                \draw[<->, rosa!80!black] (9.1, 1.5) -- (9.1, 4.5) node[midway, right] {$h$};
                \foreach \dx/\dy/\r in {1.6/0.9/3.2, -1.2/1.3/2.6, -1.9/-0.8/2.0, 0.8/-1.4/1.5} {
                        \draw[flecha, celeste!80!black, thin] (6, 3) -- (6 + \dx, 3 + \dy);
                        \fill[celeste!80!black] (6 + \dx, 3 + \dy) circle (\r pt);
                    }
                \node[nota, text width=12cm, anchor=north west] at (0, -0.8) {%
                    Para cada consulta y cabeza, la consulta predice \mPoints{}
                    desplazamientos $\Delta$ y \mPoints{} pesos $a$ (el tamaño del
                    punto). El punto muestreado es $c + \Delta \cdot (w, h)/2$ y su
                    valor, la interpolación bilineal del mapa. Se repite en los
                    \mLevels{} niveles y los pesos se normalizan sobre los
                    $\mLevels \times \mPoints$ puntos.};
            \end{tikzpicture}}
    \end{minipage}
\end{figure}

La figura de la izquierda resume una capa; la de la derecha, el paso que la hace
eficiente: cada consulta lee solo \mPoints{} puntos por nivel alrededor de su caja,
en lugar de todas las celdas de los cuatro mapas.

\subsubsection{Emparejamiento y pérdida}

Las \mQueries{} predicciones de una ventana y sus $n$ cajas anotadas se emparejan
uno a uno con el algoritmo húngaro sobre una matriz de costo
(Figura~\ref{fig:matching}):
\[
    C_{kj} = \mCostClass\, \mathcal{C}_{\text{focal}}(p_k, c_j)
    + \mCostBbox\, \lVert b_k - b_j \rVert_1
    - \mCostIou\, \mathrm{GIoU}(b_k, b_j),
\]
donde el costo de clase es la pérdida focal positiva menos la negativa que se
ahorra \citep{zhu_deformable_2021} y GIoU es el IoU generalizado, que sigue
dando señal cuando las cajas no se solapan. Sobre los pares emparejados se calcula la
pérdida
\[
    \mathcal{L} = \mWeightClass\, \mathcal{L}_{\text{focal}}
    + \mCostBbox\, \mathcal{L}_{1}
    + \mCostIou\, \mathcal{L}_{\mathrm{GIoU}},
\]
con la focal ($\alpha = \mFocalAlpha$, $\gamma = \mFocalGamma$) sobre sigmoides
independientes por clase y sin canal de ``no objeto'': la confianza de una
consulta no emparejada se empuja a cero en todas las clases. El modelo puede así
emitir una caja de alta confianza con la etiqueta equivocada, que el protocolo
mide como acierto de detección y fallo de clasificación, que es exactamente la
separación que la sección~\ref{sec:protocolo-oe2} necesita. Los sesgos de las
cabezas de clase se inicializan en $\log\frac{1-\pi}{\pi}$ con $\pi = \mPrior$,
para que al inicio casi ninguna consulta dispare.

\begin{figure}[htbp]
    \centering
    \caption{Emparejamiento húngaro}
    \label{fig:matching}
    \resizebox{\textwidth}{!}{%
    \begin{tikzpicture}[node distance=3mm and 5mm, every node/.style={font=\footnotesize}]
        \node[dato] (pred) at (0, 0.8) {\mQueries{} predicciones $(p_k, b_k)$};
        \node[dato] (gt) at (0, -0.8) {$n$ cajas anotadas $(c_j, b_j)$};
        \node[etapa] (cost) at (4.4, 0) {matriz de costo\\$\mQueries \times n$};
        \node[etapa, right=of cost] (hung) {asignación húngara\\$n$ pares};
        \node[etapa, right=of hung] (loss) {pérdida\\focal + $L_1$ + GIoU};
        \node[nota, below=1mm of hung.south, anchor=north, xshift=6mm]
        {las $\mQueries - n$ consultas restantes: focal hacia 0};
        \draw[flecha] (pred.east) -- (cost.west |- pred.east) -- (cost);
        \draw[flecha] (gt.east) -- (cost.west |- gt.east) -- (cost);
        \draw[flecha] (cost) -- (hung);
        \draw[flecha] (hung) -- (loss);
    \end{tikzpicture}}
\end{figure}

De las \mQueries{} predicciones, solo las $n$ emparejadas con una anotación reciben
las tres pérdidas; las demás aprenden a no disparar.

\subsubsection{Inferencia}

La confianza de cada consulta es el máximo de sus \nClasses{} sigmoides y la
clase, su argumento; se conservan las cajas con confianza $\geq$ \mScoreFloor{}
(el piso del protocolo) y, como para todos los modelos, a lo sumo \mMaxDet{}
por ventana. El emparejamiento uno a uno desalienta los duplicados al entrenar,
pero \mQueries{} consultas sobre una a tres llamadas dejan en inferencia cajas
repetidas y anidadas, de modo que se aplica la misma supresión de no máximos por
clase que usan las líneas base (IoU \mNmsIou{}). El presupuesto de \mQueries{}
consultas se dimensiona sobre el máximo de cajas por ventana del conjunto
(\maxPerWindow{}, sección~\ref{sec:vacias}) y no sobre la mediana
(\medianPerWindow{}).

\subsubsection{Cabeza preentrenada con cantos de aves}
\label{sec:aves}

La cabeza deformable se entrena desde cero, mientras que las líneas base parten
de un detector completo. Para acortar esa asimetría, la configuración que entra a
la comparación preentrena la misma arquitectura, con los mismos hiperparámetros,
sobre un caché aparte de cantos de aves construido con el ventaneo y el mel del
Capítulo~\ref{cap:curacion} (\file{data/processed\_extra/}), y después la afina
con primates: el afinado carga en el modelo recién construido todo tensor del
preentrenado con el mismo nombre y forma (todo salvo las cabezas de clase, que
dependen del número de clases) y sigue la receta de la propuesta.

\subsubsection{Hiperparámetros y tamaño}

La Tabla~\ref{tab:hiperparametros} reúne los hiperparámetros tal como los declara
el código.

\begin{table}[htbp]
    \centering
    \small
    \begin{threeparttable}
        \caption{Hiperparámetros de la propuesta}
        \label{tab:hiperparametros}
        \begin{tabular}{@{}L{4.4cm}L{2.6cm}L{6.4cm}@{}}
            \toprule
            \textbf{Hiperparámetro} & \textbf{Valor} & \textbf{Motivo} \\
            \midrule
            Extractor & AST & \mAstCheckpoint{}, afinado entero \\
            Paso temporal de parcheo & \mTimeStride{} & $\mTokensFreq \times \mTokensTime$ tokens, \mTokenMs{}\,ms por token \\
            Dimensión de la cabeza & \mDim{} & Cabeza pequeña sobre tokens de \mHidden{} \\
            Niveles de la pirámide & \mLevels{} & De \mCellZero{} a \mCellThree{}\,ms por celda \\
            Consultas & \mQueries{} & Muy por encima de \maxPerWindow{}, el máximo de cajas por ventana \\
            Capas del decodificador & \mDecoderLayers{} & Como Deformable DETR \\
            Cabezas / puntos por nivel & \mHeads{} / \mPoints{} & Como Deformable DETR \\
            FFN / \textit{dropout} & \mFfn{} / \mDropout{} & Como Deformable DETR \\
            Focal $\alpha$, $\gamma$ & \mFocalAlpha{}, \mFocalGamma{} & Valores de RetinaNet \\
            Costos de emparejamiento (clase, $L_1$, GIoU) & \mCostClass{}, \mCostBbox{}, \mCostIou{} & \citet{zhu_deformable_2021} \\
            Pesos de la pérdida (clase, $L_1$, GIoU) & \mWeightClass{}, \mCostBbox{}, \mCostIou{} & La clase pesa el doble al asignar que al entrenar \\
            Prior de las cabezas de clase & \mPrior{} & Casi nada dispara al inicio \\
            Compresión de entrada & \mFrontend{} & La del preentrenamiento del AST \\
            NMS por clase & IoU \mNmsIou{} & El mismo posprocesado que las líneas base \\
            \bottomrule
        \end{tabular}
    \end{threeparttable}
\end{table}

La mayoría replica los valores de Deformable DETR; los que se ajustaron a este
material son el paso temporal del parcheo, la dimensión de la cabeza, los niveles
y las consultas, cada uno con su motivo. La Tabla~\ref{tab:bloques} muestra cuánto
pesa cada bloque.

\begin{table}[htbp]
    \centering
    \small
    \begin{threeparttable}
        \caption{Parámetros por bloque}
        \label{tab:bloques}
        \begin{tabular}{@{}L{7.4cm}rr@{}}
            \toprule
            \textbf{Bloque} & \thead{Parámetros (M)} & \thead{Fracción} \\
            \midrule
            \input{\figDET/tablas/bloques} \\
            \bottomrule
        \end{tabular}
        \begin{tablenotes}[flushleft]
            \footnotesize
            \item \textit{Nota.} La compresión de entrada no aparece con parámetros
            porque el log-mel no tiene ninguno.
        \end{tablenotes}
    \end{threeparttable}
\end{table}

El reparto importa para leer la comparación: el extractor concentra el
\mParamsBackbonePct{} de los \mParams{}\,M de parámetros y la cabeza completa cabe
en \mParamsHead{}\,M. Cambiar el número de clases o ensanchar la cabeza cuesta
poco; lo que cuesta, en memoria y en tiempo, es afinar el AST.

\subsection{Líneas base}
\label{sec:baselines}

Se declaran tres líneas base y cada una controla una cosa distinta
(Figura~\ref{fig:baselines}). Las tres consumen el mismo mel del caché, pintado
como imagen con el rango en dB que registra \file{meta.json}, y devuelven cajas
en el mismo marco normalizado, de modo que comparten datos, particiones y
protocolo con la propuesta.

\begin{description}
    \item[Faster R-CNN R50-FPN v2.] El detector de dos etapas de referencia, con
        \textit{backbone} ResNet-50 y pirámide de características (FPN,
        \textit{Feature Pyramid Network}) preentrenados en COCO; el mel
        se reescala a $512 \times 1024$ píxeles y se replica en tres canales. Se
        adaptan las razones de aspecto de las anclas porque las cajas de llamadas
        van de muy anchas a muy altas, lejos de las razones de los objetos de
        COCO. Es la línea base sobre espectrogramas con precedente directo en
        bioacústica \citep{escobar-amado_computer_2024}.
    \item[YOLO26s.] Un detector denso de una etapa, pequeño
        (\cParamsYol{}\,M de parámetros), entrenado con la herramienta de
        Ultralytics desde sus pesos de COCO sobre una imagen gris de
        $512 \times 512$; es la línea base de la familia YOLO sobre
        espectrogramas \citep{gonzalez_yolo-based_2025}. Las aumentaciones de
        imagen que no tienen sentido en un espectrograma se apagan (volteos,
        mosaico, tono y saturación) y quedan un corrimiento pequeño, una escala
        y brillo.
    \item[RT-DETR-L.] Un detector de conjuntos con extractor de imagen
        preentrenado en COCO (\cParamsRtd{}\,M), también por Ultralytics. Es la
        línea base que separa las dos hipótesis que la propuesta mezcla: si una
        cabeza de tipo DETR rinde distinto por ser una cabeza DETR, o por el
        extractor de audio que la alimenta. Entró al experimento después que las
        otras dos, de forma exploratoria, y la
        sección~\ref{sec:verificacion-re21} declara la consecuencia que eso tiene
        sobre su reproducibilidad.
\end{description}

\begin{figure}[htbp]
    \centering
    \caption{Líneas base}
    \label{fig:baselines}
    \resizebox{\textwidth}{!}{%
        \begin{tikzpicture}[node distance=3mm and 6mm, every node/.style={font=\footnotesize}]
            \node[dato] (mel) {mel del caché\\$1 \times \nMels \times \nFrames$};
            \node[etapa, right=of mel, yshift=20mm] (i1) {dB $\to [0, 1]$, 3 canales\\$512 \times 1024$};
            \node[etapa, right=of mel] (i2) {dB $\to$ gris\\$512 \times 512$};
            \node[etapa, right=of mel, yshift=-20mm] (i3) {dB $\to$ gris\\$512 \times 512$};
            \node[etapa, right=of i1, fill=ambar!15, draw=ambar!80!black] (m1)
            {Faster R-CNN R50-FPN v2\\anclas adaptadas\\NMS \mNmsIou{}};
            \node[etapa, right=of i2, fill=ambar!15, draw=ambar!80!black] (m2)
            {YOLO26s (Ultralytics)\\cabeza densa\\NMS \mNmsIou{} de anidadas};
            \node[etapa, right=of i3, fill=ambar!15, draw=ambar!80!black] (m3)
            {RT-DETR-L (Ultralytics)\\cabeza de conjuntos\\NMS \mNmsIou{} de anidadas};
            \node[dato, right=of m2, xshift=6mm] (out) {cajas $c_x c_y w h$\\+ clase + confianza};
            \draw[flecha] (mel) -- (i1);
            \draw[flecha] (mel) -- (i2);
            \draw[flecha] (mel) -- (i3);
            \draw[flecha] (i1) -- (m1);
            \draw[flecha] (i2) -- (m2);
            \draw[flecha] (i3) -- (m3);
            \draw[flecha] (m1) -- (out);
            \draw[flecha] (m2) -- (out);
            \draw[flecha] (m3) -- (out);
        \end{tikzpicture}}
\end{figure}

Las tres líneas base comparten la entrada, el mel del caché pintado como imagen de
distinto tamaño, y la forma de la salida; difieren en el detector. La
Tabla~\ref{tab:configuraciones} reúne las cuatro configuraciones de la
comparación, la propuesta y las tres líneas base.

\begin{table}[htbp]
    \centering
    \small
    \begin{threeparttable}
        \caption{Configuraciones comparadas}
        \label{tab:configuraciones}
        \begin{tabular}{@{}L{3.0cm}L{3.2cm}rL{2.5cm}L{2.9cm}@{}}
            \toprule
            \textbf{Configuración} & \textbf{Extractor} & \thead{Parám.\\(M)} & \thead[l]{Preen-\\trenado} & \textbf{Qué controla} \\
            \midrule
            \input{\figDET/tablas/configuraciones} \\
            \bottomrule
        \end{tabular}
        \begin{tablenotes}[flushleft]
            \footnotesize
            \item \textit{Nota.} Parámetros según el registro de cada corrida. En la propuesta
            y en las dos de Ultralytics se entrena todo; en Faster R-CNN, el
            tallo y las dos primeras etapas del ResNet quedan congelados
            (\cTrainableFrc{}\,M entrenables de \cParamsFrc{}\,M).
        \end{tablenotes}
    \end{threeparttable}
\end{table}

Las cuatro difieren en extractor, tamaño y preentrenamiento, y la última columna
dice qué variable controla cada línea base frente a la propuesta.

\subsection{Diseño experimental}

\subsubsection{Recetas de entrenamiento}

Cada configuración se entrena con la receta de su \textit{preset}
(Tabla~\ref{tab:recetas}). Los modelos del proyecto usan AdamW con un programa
OneCycle, recorte del gradiente y las mismas aumentaciones sobre el audio:
ganancia aleatoria, mezcla de una ventana vacía a un SNR dado, corrimiento
temporal y pegado de llamadas de otras ventanas, además de un \textit{jitter} de
las cajas objetivo. Los dos modelos de Ultralytics se entrenan con su propia
herramienta y sus aumentaciones de imagen, restringidas como se dijo.
\textbf{Las recetas no están igualadas, y se declara}: cada marco usa la suya, de
modo que la comparación mide sistemas completos y no solo arquitecturas.

\begin{table}[htbp]
    \centering
    \small
    \begin{threeparttable}
        \caption{Recetas de entrenamiento}
        \label{tab:recetas}
        \begin{tabular}{@{}L{3.6cm}rrlL{2.6cm}L{3.0cm}@{}}
            \toprule
            \textbf{Configuración} & \thead{Épocas} & \thead{Lote} & \thead{lr} & \textbf{Optimizador} & \textbf{Elección de \file{best.pt}} \\
            \midrule
            \input{\figDET/tablas/recetas} \\
            \bottomrule
        \end{tabular}
    \end{threeparttable}
\end{table}

La asimetría de la última columna es una diferencia de receta, no de protocolo:
las dos corridas de Ultralytics eligen su \textit{checkpoint} con la aptitud de su
propio validador y no con la mAP@\pMatchIou{} que usan las otras dos, pero el
protocolo de comparación vuelve a medir a todos con la misma vara.

\subsubsection{Protocolo de comparación}
\label{sec:protocolo-oe2}

El protocolo vive en \file{src/evaluation/protocol.py}, se fijó antes de mirar
los resultados de prueba y la sección~\ref{sec:re23} lo aplica sin cambios
(Figura~\ref{fig:protocolo-oe2}). Consta de seis reglas, que la figura ordena
en el tiempo y que se detallan a continuación.

\begin{figure}[htbp]
    \centering
    \caption{Protocolo de comparación}
    \label{fig:protocolo-oe2}
    \resizebox{\textwidth}{!}{%
        \begin{tikzpicture}[node distance=4mm and 7mm, every node/.style={font=\footnotesize}]
            \node[etapa] (train) {entrenamiento\\receta del \textit{preset}};
            \node[dato, right=of train] (best) {best.pt};
            \node[etapa, right=of best] (dump)
            {volcado de val y test\\score $\geq$ \mScoreFloor{}, \mMaxDet{} cajas};
            \node[etapa, below=of dump] (thr)
            {val elige el umbral\\precisión $\geq$ \pMinPrecision{}, máximo \textit{recall}};
            \node[etapa, left=of thr] (test)
            {test se mide una vez\\recall, precisión, F\textsubscript{1}, IC 95\,\%};
            \node[etapa, left=of test] (per) {por eje y por clase\\IoU, AP, confusión};
            \draw[flecha] (train) -- (best);
            \draw[flecha] (best) -- (dump);
            \draw[flecha] (dump) -- (thr);
            \draw[flecha] (thr) -- (test);
            \draw[flecha] (test) -- (per);
            \node[nota, below=1mm of test] {\textit{bootstrap} por grabación};
        \end{tikzpicture}}
\end{figure}

\begin{enumerate}
    \item \textbf{Se vuelcan predicciones, no métricas.} De cada \file{best.pt}
          se guardan, por ventana de validación y de prueba, todas las cajas con
          confianza $\geq$ \mScoreFloor{}; toda cifra sale después de esos
          volcados. Cambiar una métrica o un umbral no exige volver a pasar los
          modelos por la GPU.
    \item \textbf{Mismo presupuesto para todos.} Se conservan las \mMaxDet{}
          mejores cajas por ventana: torchvision corta en 100 y Ultralytics en
          300, y sin techo común la cola de la curva precisión--\textit{recall}
          sería del marco de trabajo y no del modelo.
    \item \textbf{Acierto a IoU $\geq$ \pMatchIou{}.} Se mide si la llamada se
          \emph{encontró}, no si la caja quedó ajustada: el eje de frecuencia lo
          dibuja el anotador con criterio variable y las llamadas son cortas. A
          IoU \pStrictIou{} el corrimiento temporal tolerado es un tercio del
          ancho de la caja; a \pMatchIou{} sube a la mitad. Más abajo no: a 0,25
          una predicción de banda completa acertaría una parte de las cajas sin
          encuadrar nada. El encuadre se reporta aparte, a \pStrictIou{}. La precisión media se calcula con la
          definición y la interpolación de la curva precisión--\textit{recall} que
          fija COCO \citep{lin_2014}.
    \item \textbf{Umbral por modelo, elegido en validación.} De la rejilla
          0,01--0,99, el umbral de mayor \textit{recall} entre los que alcanzan
          una precisión de al menos \pMinPrecision{}. Un umbral fijo compartido
          mide calibración, no detección: cada cabeza reparte sus confianzas a su
          manera, y la sección~\ref{sec:carga-igual} muestra cuánto cambia la
          lectura según dónde se mida.
    \item \textbf{La prueba se mide una vez}, en ese umbral, con intervalos de
          confianza al 95\,\% por \textit{bootstrap} de \pBootstrap{} remuestreos
          \emph{de grabaciones}: las ventanas vecinas comparten la mitad del
          audio y la calidad de una grabación es de la grabación entera.
          Remuestrear ventanas daría intervalos más angostos de lo que
          corresponde.
    \item \textbf{Las clases de ventana van aparte.} Si la mediana del ancho de
          caja de una clase ocupa al menos el \pWindowWidth{} de la ventana
          (\rWindowClasses{}), detectarla es clasificar la ventana: se reportan,
          pero no entran en la mAP ni en el promedio por clase. Quedan
          \rNDetectionClasses{} clases de detección.
\end{enumerate}

\subsubsection{Métricas por eje y regla de decisión}

La Tabla~\ref{tab:metricas-eje} asigna a cada eje la pregunta que responde y la
magnitud con que se responde, en la línea que fija la
sección~\ref{sec:metricas}.

\begin{table}[htbp]
    \centering
    \small
    \caption{Métricas por eje}
    \label{tab:metricas-eje}
    \begin{tabular}{@{}L{2.2cm}L{5.0cm}L{6.2cm}@{}}
        \toprule
        \textbf{Eje} & \textbf{Pregunta} & \textbf{Magnitud} \\
        \midrule
        Detección & ¿Se encontró la llamada? & \textit{Recall} y precisión emparejando sin mirar la clase, a IoU \pMatchIou{}, al punto de operación \\
        Encuadre & ¿La caja está donde el anotador la dibujó? & IoU mediano de esos pares; \textit{recall} con clase a IoU \pStrictIou{}; AP a IoU \pLooseIou{} y \pStrictIou{} \\
        Clasificación & ¿Se etiquetó bien lo que se encontró? & Fracción de pares con la clase y con la especie correctas; \textit{recall} y AP por clase; matriz de confusión \\
        Operación & ¿Cuánto cuesta revisar? & Umbral elegido, F\textsubscript{1}, FP/TP, propuestas por hora de audio, minutos de GPU por hora de audio \\
        \bottomrule
    \end{tabular}
\end{table}

La regla de decisión se declara antes de correr el experimento y no se cambia
después:

\begin{quote}
    La métrica \textbf{primaria} es el \textit{recall} en prueba al punto de
    operación de precisión $\geq$ \pMinPrecision{}, con su intervalo de
    confianza; la \textbf{secundaria} es la mAP@\pMatchIou{}. Dos configuraciones
    cuyos intervalos se solapan \emph{no se separan} con este experimento, y se
    dice así. El modelo que pasa a RE2.4 es el de mayor \textit{recall} primario;
    si la primaria no separa, decide la secundaria, y en caso de empate también
    en ella, el más barato en inferencia.
\end{quote}

La comparación pone a prueba una hipótesis de diseño: que un extractor
preentrenado en audio con una cabeza de conjuntos supera a los detectores
preentrenados en las imágenes de COCO cuando todos se entrenan sobre el mismo
mel. Con la regla anterior, se confirma si la propuesta queda primera en la
métrica primaria o, si esta no separa, en la secundaria. La
sección~\ref{sec:seleccion-modelo} da el veredicto.

\subsection{Verificación de RE2.1}
\label{sec:verificacion-re21}

El indicador pide que se especifique la arquitectura implementada, con sus
hiperparámetros, y que se defina la comparación con las líneas base declaradas.
La Tabla~\ref{tab:verificacion-re21} contrasta cada elemento con el código.

\begin{table}[htbp]
    \centering
    \small
    \begin{threeparttable}
        \caption{Verificación de RE2.1}
        \label{tab:verificacion-re21}
        \begin{tabular}{@{}L{3.6cm}L{2.2cm}L{7.6cm}@{}}
            \toprule
            \textbf{Elemento del IOV} & \textbf{Valor} & \textbf{Dónde se comprueba} \\
            \midrule
            Arquitectura AST-Deformable-DETR implementada & sí & \file{src/models/deformable\_detr.py} y \file{backbone.py}; la Figura~\ref{fig:arquitectura} es su recorrido \\
            Todos sus hiperparámetros & \mParams{}\,M de parámetros & Tablas~\ref{tab:hiperparametros} y~\ref{tab:bloques}, generadas desde el código y desde el \file{best.pt}: dimensión \mDim{}, \mQueries{} consultas, \mLevels{} niveles \\
            Comparación definida antes de entrenar: configuraciones & 3 líneas base & Sección~\ref{sec:baselines} y Tabla~\ref{tab:configuraciones} \\
            Partición, métricas y regla de selección & fijadas antes de entrenar & Protocolo de la sección~\ref{sec:protocolo-oe2}, que aplica el de la sección~\ref{sec:metodos} \\
            \bottomrule
        \end{tabular}
    \end{threeparttable}
\end{table}

\rotulo{Por qué \mQueries{} consultas y \mLevels{} niveles.} El prototipo con
que se planificó el proyecto usaba 64 consultas y 3 niveles, y la implementación
final usa \mQueries{} y \mLevels{}. Ninguno de los dos cambios altera la
arquitectura:

\begin{itemize}[nosep]
    \item \textbf{De 64 a \mQueries{} consultas.} El presupuesto de consultas
          acota cuántos eventos puede emitir el modelo en una ventana. El máximo
          observado en el conjunto es \maxPerWindow{} cajas por ventana, de modo
          que 64 ya bastaba; subirlo a \mQueries{} añade
          $\mQueries \times \mDim$ parámetros de incrustación (una fracción
          despreciable de los \mParams{}\,M) y deja margen para escenas más
          densas que las de este corpus. El costo es cuadrático en la
          autoatención entre consultas, que sobre \mQueries{} elementos es
          irrelevante frente al extractor.
    \item \textbf{De 3 a \mLevels{} niveles.} El nivel que se añadió es el
          $\tfrac{1}{2}\times$, el más grueso (\mCellThree{}\,ms por celda), y la
          razón está en la Figura~\ref{fig:duraciones}: con tres niveles la
          resolución más gruesa habría sido \mCellTwo{}\,ms, por debajo de la
          duración mediana de buena parte de las clases, y el decodificador
          habría tenido que componer los eventos largos a partir de celdas finas.
\end{itemize}

\rotulo{Una cuarta configuración, incorporada después.} RT-DETR-L se lanzó de
forma exploratoria con la herramienta de Ultralytics, fuera del registro de
\textit{presets}, y su receta se recuperó del \file{config.json} de la corrida
para dejarla en \file{src/train.py} como \code{--arch rtdetr}. Comparte el flujo
de YOLO26s salvo en dos puntos, que son los que la corrida original fijó a mano:
el optimizador, AdamW con tasa inicial \(10^{-4}\), porque con el que elige
Ultralytics por su cuenta la configuración no converge; y el lote, \(16\) en
lugar de \(32\), porque el decodificador de conjuntos no entra en memoria con el
lote de YOLO. Su \textit{checkpoint} se carga en modo estricto y reproduce su
volcado caja por caja, de modo que las cifras que aquí se reportan proceden del
mismo código que se publica.

% ===================================================================
\section{Pipeline reproducible de entrenamiento, evaluación e inferencia (RE2.2)}
\label{sec:re22}
% ===================================================================

\rotulo{Resultado.} Un \textit{pipeline} de ocho puntos de entrada que va de las
tablas de Raven a las tablas de Raven: prepara los datos, entrena desde una
configuración, evalúa un \textit{checkpoint} y ejecuta inferencia sobre un WAV,
con las dependencias y los comandos declarados.

\rotulo{Cómo se alcanzó.} Con los procedimientos de la
sección~\ref{sec:procedimientos}: material primario intocado, experimentos
registrados y entorno declarado. Cada etapa es un \textit{script} de \file{src/}
con su propio \code{--help}; ninguna modifica los artefactos de la anterior, de modo que
se puede volver a ejecutar cualquiera desde su entrada
(Figura~\ref{fig:pipeline}). Las etapas de datos son las del
Capítulo~\ref{cap:curacion}; las de modelo, las de la sección anterior.

\rotulo{Cómo reproducirlo.} Con el entorno que fija \file{uv.lock}, se ejecutan las
cuatro etapas en orden, cada una desde su \textit{script}. La semilla única, la
configuración que viaja con cada \textit{checkpoint} y los controles que
impiden mezclar versiones del caché hacen que el recorrido se repita; la
sección~\ref{sec:reproducibilidad} detalla qué queda fijado y qué no.

\rotulo{Medio de verificación.} Los \textit{scripts} de \file{src/},
\file{pyproject.toml} y \file{uv.lock}, publicados en el repositorio del proyecto
(\url{\rRepoUrl}), y esta sección como manual técnico; la
sección~\ref{sec:verificacion-re22} recorre el indicador etapa por etapa. El
indicador nombra los módulos de \file{src/pipelines/}, la organización del
prototipo; en la implementación final esos módulos son los paquetes
\file{src/data/}, \file{src/models/}, \file{src/training/},
\file{src/evaluation/} y \file{src/inference/}, con un \textit{script} por
etapa. El registro de entrenamiento de cada corrida está en la carpeta RE2.2 de los
medios de verificación (Anexo~\ref{anx:medios}).

\begin{figure}[htbp]
    \centering
    \caption{Pipeline del proyecto}
    \label{fig:pipeline}
    \resizebox{0.86\textwidth}{!}{%
        \begin{tikzpicture}[node distance=3.4mm and 7mm,
                caja/.append style={font=\scriptsize, inner sep=3pt},
                dato/.append style={font=\scriptsize\ttfamily}]
            \node[dato] (raw) {data/raw/};
            \node[etapa, below=of raw] (pa) {prepare\_annotations.py};
            \node[dato, below=of pa] (cleaned) {data/cleaned/};
            \node[etapa, below=of cleaned] (pd) {prepare\_data.py};
            \node[dato, below=of pd] (proc) {data/processed/};
            \node[etapa, below=of proc] (train) {train.py --arch \{detr, detr\_resnet, frcnn, yolo\}};
            \node[etapa, right=of proc] (yolo) {export\_yolo.py};
            \node[dato, right=of yolo] (yolodir) {data/yolo/};
            \node[dato, below=of train] (run)
            {runs/<name>/: best.pt, last.pt, config.json, metrics.jsonl, train.log};
            \node[etapa, below=of run] (dump) {dump\_predictions.py};
            \node[dato, below=of dump] (dumps) {runs/<name>/<name>\_val|test\_predictions.pt};
            \node[etapa, below=of dumps] (cmp) {compare\_models.py};
            \node[dato, below=of cmp] (cmpout) {runs/comparacion/*.txt, *.json};
            \node[etapa, below=of cmpout] (detect) {detect.py};
            \node[dato, right=of detect] (op) {runs/<name>/operating\_point.json};
            \node[dato, below=of detect] (raven) {<audio>.detections.txt (tabla de Raven)};
            \foreach \a/\b in {raw/pa, pa/cleaned, cleaned/pd, pd/proc, proc/train, train/run,
                    run/dump, dump/dumps, dumps/cmp, cmp/cmpout, cmpout/detect, detect/raven}
            \draw[flecha] (\a) -- (\b);
            \draw[flecha] (proc) -- (yolo);
            \draw[flecha] (yolo) -- (yolodir);
            \draw[flecha] (yolodir.south) |- (train.east);
            \draw[flecha] (cmp.east) -| (op.north);
            \draw[flecha] (op) -- (detect);
            \draw[flecha] (run.east) -- ++(5mm, 0) |- ([yshift=1.5mm]detect.east);
            \node[nota, left=2mm of pa.west, anchor=east] {RE1.1};
            \node[nota, left=2mm of pd.west, anchor=east] {RE1.2};
            \node[nota, left=2mm of train.west, anchor=east] {RE2.1};
            \node[nota, left=2mm of cmp.west, anchor=east] {RE2.3};
            \node[nota, left=2mm of detect.west, anchor=east] {RE2.4};
            \node[nota, right=1mm of yolodir.east, anchor=west, align=left]
            {solo Ultralytics:\\PNG y etiquetas};
        \end{tikzpicture}}
\end{figure}

La figura se lee de arriba abajo: las dos primeras etapas preparan los datos
(RE1.1 y RE1.2), la tercera entrena (RE2.1), la cuarta evalúa (RE2.3) y la última
escribe la tabla de Raven (RE2.4). El desvío de la derecha es la exportación que
solo necesitan los modelos de Ultralytics.

\subsection{Dependencias declaradas}

El entorno está declarado y bloqueado, que es la condición para que el
\textit{pipeline} sea reproducible y no solo ejecutable en una máquina.
\file{pyproject.toml} fija Python $\geq$~\rPython{} y las dependencias base
(PyTorch, torchvision, torchaudio, pandas y soundfile) y separa en dos
extras lo que solo necesita una familia de arquitecturas: \code{detr}
(transformers, timm) para las dos basadas en DETR y \code{yolo} (ultralytics,
opencv) para las dos de Ultralytics. Esa separación tiene una consecuencia
práctica: una instalación de inferencia no arrastra las librerías de
entrenamiento. \file{uv.lock} fija las versiones exactas de todo el árbol.

Los pesos preentrenados (el AST de AudioSet, el Deformable DETR de COCO) se
resuelven en el primer uso y quedan en caché local, de modo que las corridas
posteriores no dependen de la red. El entrenamiento se hizo sobre una GPU de
24\,GB; la inferencia no necesita GPU.

\subsection{Las cuatro etapas y su interfaz}

El indicador de este resultado exige que los comandos de ejecución queden
declarados. El Listado~\ref{lst:pipeline} los recoge: es la interfaz del
\textit{pipeline}, no un instructivo, y cada línea corresponde a una caja de la
Figura~\ref{fig:pipeline}.

\begin{lstlisting}[language=bash, caption={Interfaz del pipeline, una línea por etapa}, label={lst:pipeline}]
# 1. datos (capitulo 4)
prepare_annotations.py [--force]     raw/ -> cleaned/
prepare_data.py        [--force]     cleaned/ -> processed/
export_yolo.py                       processed/ -> yolo/ (solo Ultralytics)

# 2. entrenamiento: --arch elige el preset; --hp y --cfg lo pisan
train.py --arch detr  --name detr_t10_logmel --hp frontend=logmel
train.py --arch frcnn --name frcnn_v3
train.py --arch yolo  --name yolo26s_v3 --hp model=yolo26s
train.py --arch detr  --limit 64 --cfg epochs=1      # prueba de humo
notebooks/pretrain_birds_train.py --name detr_birds_v2
notebooks/pretrain_birds_train.py --finetune <aves>/best.pt --name <corrida>

# 3. evaluación: volcar primero, medir después
dump_predictions.py --run <corrida> [--device cuda:1]
compare_models.py runs/*/*_predictions.pt --output <salida>
./evaluation.sh [corridas...]        las dos cosas, solo lo desactualizado

# 4. inferencia
detect.py <audio o carpeta> --model runs/<corrida>/best.pt
\end{lstlisting}

Una corrida se identifica por su \textit{preset} (la arquitectura del registro
\code{models.\allowbreak registry.\allowbreak ARCHITECTURES} con sus
hiperparámetros y su \code{TrainConfig}) y por lo que se le pisa encima.
Cualquier valor que se cambie respecto del \textit{preset} queda escrito en el
nombre de la corrida, en su \file{config.json} y dentro del propio
\textit{checkpoint}: es lo que permite, meses después, saber con qué se entrenó
un archivo de pesos sin depender de ninguna anotación externa.

El ciclo de entrenamiento (Figura~\ref{fig:trainer}) es común a las
arquitecturas del proyecto; las de Ultralytics delegan en su herramienta con las
aumentaciones restringidas de la sección anterior y al terminar exportan su mejor
época al formato de \textit{checkpoint} del proyecto, de modo que aguas abajo las
cuatro se tratan igual.

\begin{figure}[htbp]
    \centering
    \caption{Una época del entrenador}
    \label{fig:trainer}
    \resizebox{0.88\textwidth}{!}{%
        \begin{tikzpicture}[node distance=3.4mm and 8mm,
                caja/.append style={font=\scriptsize, inner sep=3pt},
                dato/.append style={font=\scriptsize\ttfamily}]
            \node[dato] (data) {train.pt (SpectrogramDataset)};
            \node[etapa, below=of data] (aug)
            {aumentación: ganancia, fondo, corrimiento, pegado; \textit{jitter} de cajas};
            \node[etapa, below=of aug] (step)
            {paso: pérdida $\to$ retropropagación $\to$ recorte $\to$ AdamW};
            \node[etapa, below=of step] (val)
            {validación: \code{detect} sobre val.pt $\to$ recall, precisión, mAP@\pMatchIou{}};
            \node[dato, below=of val] (log) {metrics.jsonl, train.log};
            \node[decision, below=of log, xshift=-34mm] (best) {¿mejor mAP@\pMatchIou{}?};
            \node[dato, below=of best] (bestpt) {best.pt};
            \node[dato, below=of log, xshift=34mm] (lastpt) {last.pt (con optimizador y programa)};
            \draw[flecha] (data) -- (aug);
            \draw[flecha] (aug) -- (step);
            \draw[flecha] (step) -- (val);
            \draw[flecha] (val) -- (log);
            \draw[flecha] (log) -- (best);
            \draw[flecha] (log) -- (lastpt);
            \draw[flecha] (best) -- node[left, nota] {sí} (bestpt);
            \draw[flecha, dashed] (val.east) -- ++(6mm, 0) |- (aug.east)
            node[pos=0.25, right, nota] {siguiente época};
            \node[nota, right=2mm of lastpt.east, anchor=west, align=left]
            {reanuda una corrida\\interrumpida};
        \end{tikzpicture}}
\end{figure}

Cada época aumenta los datos, entrena, valida y registra. El \file{best.pt} solo
se reemplaza cuando mejora la mAP@\pMatchIou{} de validación, y el \file{last.pt}
guarda además el estado del optimizador para reanudar una corrida interrumpida.

\subsection{Artefactos de una corrida y autonomía del \textit{checkpoint}}

Una corrida deja su \file{config.json} con la arquitectura, los hiperparámetros y
el \file{meta.json} del caché con el que se entrenó; su \file{train.log} y su
\file{metrics.jsonl} con una línea por época; el \file{best.pt} de la mejor época
y el \file{last.pt} con optimizador y programa para reanudar; después de la
evaluación, los dos volcados y su \file{operating\_point.json}.

Un \textit{checkpoint} es autocontenido (Tabla~\ref{tab:checkpoint}):
\code{models.\allowbreak registry.\allowbreak load\_\allowbreak checkpoint(\allowbreak path)}
reconstruye el modelo a partir de
\code{architecture}, \code{hparams} y el número de clases de \code{labels}, carga
el \code{state\_dict} en modo estricto (ninguna clave puede faltar ni
sobrar) y devuelve el modelo junto con las etiquetas, la configuración y el
umbral de operación si existe. Es lo que abren \file{detect.py} y
\file{dump\_predictions.py}, y es la propiedad que RE2.4 verifica.

\begin{table}[htbp]
    \centering
    \small
    \caption{Claves de un checkpoint}
    \label{tab:checkpoint}
    \begin{tabular}{@{}lL{4.6cm}L{6.4cm}@{}}
        \toprule
        \textbf{Clave} & \textbf{Contenido} & \textbf{Para qué} \\
        \midrule
        \input{\figDET/tablas/checkpoint} \\
        \bottomrule
    \end{tabular}
\end{table}

Las claves de arquitectura, hiperparámetros y etiquetas bastan para reconstruir el
modelo sin el código de entrenamiento ni el caché, y eso es lo que permite
publicar un único archivo como modelo final.

\subsection{Costo de las corridas}

La Tabla~\ref{tab:costo} registra las cuatro corridas de la comparación: cuántas
épocas corrieron de las planificadas, en cuál quedó el \textit{checkpoint}
elegido y cuánto costó cada una. El total es de \cHoursTotal{} horas de GPU, con
un reparto muy desigual: \cHoursYol{}\,h la más barata y \cHoursDet{}\,h la
propuesta, que es la única que afina un extractor de \cParamsDet{}\,M de
parámetros.

\begin{table}[htbp]
    \centering
    \small
    \begin{threeparttable}
        \caption{Costo de las corridas}
        \label{tab:costo}
        \begin{tabular}{@{}L{4.6cm}rrrr@{}}
            \toprule
            \textbf{Configuración} & \thead{Épocas} & \thead{Época\\elegida} & \thead{Minutos\\por época} & \thead{Horas\\de GPU} \\
            \midrule
            \input{\figDET/tablas/costo} \\
            \bottomrule
        \end{tabular}
        \begin{tablenotes}[flushleft]
            \footnotesize
            \item \textit{Nota.} Horas de GPU: diferencia entre la primera y la última marca de
            tiempo del registro de la corrida, de modo que incluye la validación
            de cada época. La GPU se compartió con otros procesos: son cotas
            superiores y no son estrictamente comparables entre corridas. Las dos
            corridas de Ultralytics pararon antes de las 30 épocas por su
            criterio de paciencia. La corrida de la propuesta cuenta solo el
            afinado con primates, no el preentrenamiento con aves.
        \end{tablenotes}
    \end{threeparttable}
\end{table}

La Figura~\ref{fig:curvas} muestra las curvas de validación. Las dos corridas del
proyecto se miden con la mAP@\pMatchIou{} del protocolo y las dos de Ultralytics
con la mAP@0,5 de su propio validador: no son la misma magnitud y por eso van en
paneles separados.

\begin{figure}[htbp]
    \centering
    \caption{Curvas de entrenamiento}
    \label{fig:curvas}
    \begin{tikzpicture}
        \begin{groupplot}[
                group style={group size=2 by 1, horizontal sep=1.6cm},
                width=0.48\textwidth, height=5.6cm,
                xlabel={época}, xmin=0, xmax=31,
                axis line style={black!50}, tick style={black!50},
                grid=major, grid style={black!12},
                legend style={font=\scriptsize, draw=black!30, fill=white,
                        at={(0.98,0.03)}, anchor=south east},
            ]
            \nextgroupplot[ylabel={mAP@\pMatchIou{} en validación}, ymin=0.3, ymax=0.68]
            \addplot[azul, thick, mark=none] table[x=epoca, y=map] {\figDET/data/curva_det.dat};
            \addlegendentry{AST-Def-DETR}
            \addplot[ambar!85!black, thick, mark=none] table[x=epoca, y=map] {\figDET/data/curva_frc.dat};
            \addlegendentry{Faster R-CNN}
            \nextgroupplot[ylabel={mAP@0,5 de Ultralytics}, ymin=0.0, ymax=0.68]
            \addplot[verde!70!black, thick, mark=none] table[x=epoca, y=map] {\figDET/data/curva_yol.dat};
            \addlegendentry{YOLO26s}
            \addplot[rosa!80!black, thick, mark=none] table[x=epoca, y=map] {\figDET/data/curva_rtd.dat};
            \addlegendentry{RT-DETR-L}
        \end{groupplot}
    \end{tikzpicture}
\end{figure}

Las cuatro curvas tienen la misma forma: la mayor parte de la mejora ocurre en el
primer tercio del programa y después la curva se aplana. La propuesta alcanza su
máximo en la época \cBestEpochDet{} de \cEpochsDet{} y Faster R-CNN en la
\cBestEpochFrc{} de \cEpochsFrc{}; prolongar el programa no habría cambiado la
elección.

\subsection{Reproducibilidad}
\label{sec:reproducibilidad}

\begin{itemize}[nosep]
    \item \textbf{Semillas.} Una semilla única (\pSeed{}) fija la partición, las
          ventanas vacías, el orden de los lotes, las aumentaciones y el
          \textit{bootstrap} del protocolo; las corridas de Ultralytics reciben la
          misma y corren en modo determinista.
    \item \textbf{Trazabilidad.} Cada \textit{checkpoint} lleva dentro el
          \file{meta.json} del caché con el que se entrenó y sus hiperparámetros;
          cada corrida, su \file{config.json} y su registro con marcas de tiempo.
          Una cifra de este capítulo se rastrea a un \textit{commit}
          (\rCommit{}), un \file{meta.json} y un \file{best.pt}.
    \item \textbf{Fallo temprano.} \code{prepare\_*} se niegan a pisar salidas,
          el entrenamiento de Ultralytics se niega a entrenar sobre un export de
          otra versión del caché, \code{compare\_models.py} se niega a comparar
          volcados que no evaluaron las mismas ventanas, y \code{load\_checkpoint}
          carga en modo estricto.
    \item \textbf{Pruebas de humo.} \code{--limit N --cfg epochs=1 workers=0}
          recorre el ciclo completo en minutos; es la prueba que precede a cada
          corrida larga.
    \item \textbf{Lo que no está fijado.} El orden de las operaciones en GPU no es
          determinista en todos los núcleos de cuDNN: dos corridas con la misma
          semilla dan curvas parecidas, no idénticas. El costo de una corrida
          depende de la carga de la GPU compartida.
\end{itemize}

\subsection{Verificación de RE2.2}
\label{sec:verificacion-re22}

El indicador enumera cuatro capacidades y dos requisitos de declaración. La
Tabla~\ref{tab:verificacion-re22} las recorre con la evidencia de cada una.

\begin{table}[htbp]
    \centering
    \small
    \begin{threeparttable}
        \caption{Verificación de RE2.2}
        \label{tab:verificacion-re22}
        \begin{tabular}{@{}L{3.3cm}L{4.6cm}L{5.5cm}@{}}
            \toprule
            \textbf{Capacidad} & \textbf{Comando} & \textbf{Evidencia} \\
            \midrule
            Genera las particiones procesadas & \file{prepare\_annotations.py}, \file{prepare\_data.py} & \file{data/processed/} con los tres \file{.pt}, \file{labels.json} y \file{meta.json} (sección~\ref{sec:verificacion-re12}) \\
            Entrena desde un conjunto de configuración & \file{train.py --arch} & Las cuatro corridas de la Tabla~\ref{tab:costo}, cada una con su \file{config.json} \\
            Evalúa un \textit{checkpoint} & \file{dump\_predictions.py}, \file{compare\_models.py}, \file{evaluation.sh} & \file{runs/comparacion/} con el umbral, las métricas y el reparto por clase de cada corrida, publicados en la carpeta RE2.3 del Anexo~\ref{anx:medios} \\
            Ejecuta inferencia sobre un WAV & \file{detect.py} & La demostración de la sección~\ref{sec:re24}, con su transcripción \\
            Dependencias declaradas & \file{pyproject.toml}, \file{uv.lock} & Python $\geq$ \rPython{}, extras \code{detr} y \code{yolo} \\
            Comandos declarados & Listado~\ref{lst:pipeline} & Los ocho puntos de entrada de la Figura~\ref{fig:pipeline} \\
            \bottomrule
        \end{tabular}
    \end{threeparttable}
\end{table}

Las seis filas tienen su evidencia en el documento o en el repositorio. La única
que depende de material no publicado es la primera, porque el caché se construye
con las grabaciones del equipo de investigación.

% ===================================================================
\section{Informe comparativo de resultados (RE2.3)}
\label{sec:re23}
\label{sec:resultados-deteccion}
% ===================================================================

\rotulo{Resultado.} La comparación de las cuatro configuraciones sobre la misma
partición de prueba, con el mismo protocolo, separada en los tres ejes (
detección, encuadre y clasificación) y con la selección del modelo justificada
por la regla declarada en la sección~\ref{sec:protocolo-oe2}.

\rotulo{Cómo se alcanzó.} Aplicando la Fase~5 de CRISP-DM
(sección~\ref{sec:metricas}), con sus métricas y su criterio de emparejamiento.
\file{evaluation.sh} vuelca las predicciones crudas de cada \file{best.pt} sobre
validación y prueba y aplica el protocolo sin volver a pasar los modelos por la
GPU. Las cifras de validación
eligen el umbral de cada modelo; las de prueba se miden una sola vez con ese
umbral.

\rotulo{Cómo reproducirlo.} Con los cuatro \file{best.pt} y el caché de
\file{data/processed/}, \file{evaluation.sh} vuelve a volcar las predicciones y
\file{compare\_models.py} recalcula todas las métricas. El umbral se elige con una
regla fija y el \textit{bootstrap} usa la semilla \pSeed{}, de modo que las
cifras salen idénticas. Solo se publica el \textit{checkpoint} del modelo final
(RE2.4); los otros tres se obtienen entrenando con sus recetas (RE2.1).

\rotulo{Medio de verificación.} Esta sección, que es el informe de
experimentación, y los archivos de métricas que produce \file{compare\_models.py}
(\file{comparacion\_\allowbreak modelos\_\allowbreak v3.txt} y su \file{.json}),
de los que se generan todas las tablas y figuras de la sección; la
sección~\ref{sec:verificacion-re23} contrasta el indicador magnitud por magnitud.
Los dos archivos de métricas, el punto de operación de cada corrida y la matriz de
confusión están en la carpeta RE2.3 de los medios de verificación
(Anexo~\ref{anx:medios}).

\subsection{Lo que se compara}

Las cuatro configuraciones de la Tabla~\ref{tab:configuraciones} se entrenaron
sobre el mismo conjunto derivado de \nClasses{} clases, con la misma partición
por grabación, y se miden sobre las mismas \nWindowsTest{} ventanas de prueba
(\nBoxesTest{} cajas anotadas, \rHoursTest{} horas de audio). El protocolo
comprueba esa igualdad antes de comparar nada y se niega a seguir si dos volcados
no evaluaron las mismas ventanas.

\subsection{Desempeño al punto de operación}

La Tabla~\ref{tab:resultados-op} reporta el umbral que la validación eligió para
cada modelo y lo que ese umbral rinde en prueba.

\begin{table}[htbp]
    \centering
    \small
    \begin{threeparttable}
        \caption{Resultados al punto de operación}
        \label{tab:resultados-op}
        \setlength{\tabcolsep}{3pt}
        \scriptsize
        \begin{tabular}{@{}L{2.9cm}ccccccccc@{}}
            \toprule
            \textbf{Configuración} & \thead{Umbral} & \thead{Recall\\(val)} & \thead{Recall\\(test)} & \thead{IC 95\,\%} & \thead{Prec.} & \thead{F\textsubscript{1}} & \thead{mAP\\@\pLooseIou{}} & \thead{mAP\\@\pMatchIou{}} & \thead{mAP\\@\pStrictIou{}} \\
            \midrule
            \input{\figDET/tablas/global} \\
            \bottomrule
        \end{tabular}
        \begin{tablenotes}[flushleft]
            \footnotesize
            \item \textit{Nota.} El \textit{recall} y la precisión de esta tabla exigen la
            etiqueta correcta; la Tabla~\ref{tab:ejes} los separa del acierto de
            detección. Las mAP se promedian sobre las
            \rNDetectionClasses{} clases de detección, sin las dos de ventana.
            \item El intervalo es de \pBootstrap{} remuestreos de grabaciones, no
            de ventanas.
        \end{tablenotes}
    \end{threeparttable}
\end{table}

Lo primero que hay que registrar, porque condiciona todo lo que sigue, es que
\textbf{la métrica primaria no separa a ninguna pareja de configuraciones}: los
cuatro intervalos de confianza se solapan, el \textit{recall} de prueba va de
\rRecallFrc{} a \rRecallYol{} (una diferencia de \rRecallSpread{}) y las
cuatro F\textsubscript{1} caben en tres milésimas. Por la regla declarada en la
sección~\ref{sec:protocolo-oe2}, eso no se redondea a un ganador: se dice que el
experimento no las distingue en esa magnitud, y se pasa a la secundaria.

La métrica secundaria sí ordena, y con un margen de \rMapSpread{}:

\begin{center}
    YOLO26s (\rMapYol{}) $>$ RT-DETR-L (\rMapRtd{}) $>$
    AST-Def-DETR (\rMapDet{}) $\approx$ Faster R-CNN (\rMapFrc{})
\end{center}

La mAP es libre de umbral: integra toda la curva precisión--\textit{recall} en
lugar de leer un punto, de modo que no depende de dónde cada cabeza puso su
escala de confianza. Que ordene donde el punto de operación no ordena significa
que las cuatro configuraciones alcanzan una cobertura equivalente al presupuesto
de revisión que fija \pMinPrecision{} de precisión, pero no con la misma calidad
de curva: dos de ellas sostienen mejor la precisión cuando se les pide más
cobertura. El mismo orden se mantiene a IoU \pLooseIou{} y a IoU \pStrictIou{}, y
la distancia crece al apretar el encuadre (\rMapStrictYol{} contra
\rMapStrictDet{}), lo que apunta a que la diferencia está en cómo encuadran y
no en cuánto encuentran.

\rotulo{Los umbrales elegidos son muy distintos entre sí.} \rThresholdYol{} para
YOLO26s y \rThresholdFrc{} para Faster R-CNN, con la propuesta en
\rThresholdDet{} y RT-DETR-L en \rThresholdRtd{}. Esto es exactamente lo que la
cuarta regla del protocolo anticipaba: un umbral fijo compartido (por ejemplo
el 0,5 de fábrica) habría situado a cada modelo en un punto distinto de su
curva y habría medido calibración tanto como detección. La
sección~\ref{sec:carga-igual} desarrolla la consecuencia.

\subsection{Detección, encuadre y clasificación}
\label{sec:ejes}

La Tabla~\ref{tab:ejes} separa los tres ejes. Es la tabla más importante del
capítulo, porque es la que dice dónde está el error.

\begin{table}[htbp]
    \centering
    \small
    \begin{threeparttable}
        \caption{Detección, encuadre y clasificación}
        \label{tab:ejes}
        \begin{tabular}{@{}L{3.6cm}cccccc@{}}
            \toprule
                                   & \multicolumn{2}{c}{\textbf{Detección}} & \multicolumn{2}{c}{\textbf{Encuadre}} & \multicolumn{2}{c}{\textbf{Clasificación}} \\
            \cmidrule(lr){2-3}\cmidrule(lr){4-5}\cmidrule(lr){6-7}
            \textbf{Configuración} & \thead{Recall} & \thead{Prec.} & \thead{IoU\\mediano} & \thead{Recall\\@\pStrictIou{}} & \thead{Clase\\correcta} & \thead{Especie\\correcta} \\
            \midrule
            \input{\figDET/tablas/ejes} \\
            \bottomrule
        \end{tabular}
        \begin{tablenotes}[flushleft]
            \footnotesize
            \item \textit{Nota.} Detección: emparejando sin mirar la clase. Encuadre: IoU
            mediano de esos mismos pares, y \textit{recall} exigiendo la clase a
            IoU \pStrictIou{}. Clasificación: de lo que se encontró, qué fracción
            lleva la clase y la especie anotadas.
        \end{tablenotes}
    \end{threeparttable}
\end{table}

\rotulo{La clasificación no es el problema.} Entre \rClassOkMin{} y
\rClassOkMax{} de los pares emparejados llevan la clase correcta, y al menos
\rSpeciesOkMin{} la especie correcta, en las cuatro configuraciones. Distinguir
\nClasses{} pares de especie y tipo de llamada, con \classImbalance{} veces más
cajas de entrenamiento en la clase más frecuente que en la menos frecuente, cuesta
menos de dos puntos una vez que el evento está localizado. La distancia entre el \textit{recall} de
detección de la Tabla~\ref{tab:ejes}, que empareja sin mirar la clase, y el de la
Tabla~\ref{tab:resultados-op}, que sí exige la etiqueta, mide ese costo: para el
modelo seleccionado es de poco más de un punto.

\rotulo{El encuadre sí separa.} El IoU mediano va de \rIouDet{} a \rIouRtd{}, y la
fracción de pares que superan IoU \pStrictIou{} (el umbral estricto) va de
\rIouStrictShareDet{} a \rIouStrictShareYol{}. La Figura~\ref{fig:iou} muestra la
distribución completa: las cuatro tienen la misma forma, una moda ancha entre 0,7
y 0,9, pero las dos configuraciones de Ultralytics la desplazan a la derecha.
Esta es la diferencia que la mAP@\pStrictIou{} amplifica.

\begin{figure}[htbp]
    \centering
    \caption{Distribución del IoU}
    \label{fig:iou}
    \begin{tikzpicture}
        \begin{axis}[
                width=0.92\textwidth, height=6.4cm,
                xlabel={IoU del par emparejado}, ylabel={fracción de los pares},
                xmin=0.3, xmax=1.0, ymin=0,
                axis line style={black!50}, tick style={black!50},
                grid=major, grid style={black!12},
                legend style={font=\scriptsize, draw=black!30, fill=white,
                        at={(0.02,0.98)}, anchor=north west},
            ]
            \addplot[azul, thick, mark=none] table[x=iou, y=Det] {\figDET/data/iou.dat};
            \addlegendentry{AST-Def-DETR}
            \addplot[ambar!85!black, thick, mark=none] table[x=iou, y=Frc] {\figDET/data/iou.dat};
            \addlegendentry{Faster R-CNN}
            \addplot[verde!70!black, thick, mark=none] table[x=iou, y=Yol] {\figDET/data/iou.dat};
            \addlegendentry{YOLO26s}
            \addplot[rosa!80!black, thick, mark=none] table[x=iou, y=Rtd] {\figDET/data/iou.dat};
            \addlegendentry{RT-DETR-L}
            \draw[black!55, dashed] (axis cs:0.5,0) -- (axis cs:0.5,0.30)
            node[above, font=\scriptsize, text=black!65] {IoU \pStrictIou{}};
        \end{axis}
    \end{tikzpicture}
\end{figure}

La matriz de confusión del modelo seleccionado (Figura~\ref{fig:confusion})
confirma la lectura por el otro lado.

\begin{figure}[htbp]
    \centering
    \caption{Matriz de confusión}
    \label{fig:confusion}
    \includegraphics[width=0.96\textwidth]{\figDET/matriz_confusion.png}
\end{figure}

Sobre \rConfusionPairs{} pares emparejados, solo el \rConfusionOff{} de la masa cae
fuera de la diagonal, y el \rSameSpeciesShare{} de esos errores es entre tipos de
llamada de la \emph{misma} especie: el error que anticipaba la geometría del
conjunto, que separa especies pero no tipos de llamada dentro de una especie
(sección~\ref{sec:geometria}). La confusión más frecuente, \rTopConfusion{}, suma
\rTopConfusionN{} pares, el \rTopConfusionShare{} de todos los errores de
etiqueta.

\subsection{Cobertura frente a carga de revisión}
\label{sec:carga-igual}

La confianza que emite un detector no es una probabilidad comparable entre
arquitecturas: depende de la función de pérdida, del número de propuestas que la
cabeza produce y de cómo se normaliza la puntuación. Los cuatro umbrales elegidos
(de \rThresholdYol{} a \rThresholdFrc{}) lo muestran sin ambigüedad.

La magnitud que sí es comparable, y además es la que interesa a OE3, es
\textbf{cuántos eventos recupera cada modelo cuando a todos se les permite
    proponer la misma cantidad de cajas por hora de audio}. La
Figura~\ref{fig:carga} barre el umbral de cada configuración y lee su curva a la
misma abscisa; la Tabla~\ref{tab:carga-igual} la tabula en seis presupuestos.

\begin{figure}[htbp]
    \centering
    \caption{Cobertura frente a carga}
    \label{fig:carga}
    \begin{tikzpicture}
        \begin{axis}[
                width=0.92\textwidth, height=7.2cm,
                xlabel={propuestas por hora de audio}, ylabel={cobertura (\textit{recall})},
                xmin=400, xmax=9000, ymin=0.2, ymax=0.98,
                scaled x ticks=false, xtick={1000,2000,...,9000},
                axis line style={black!50}, tick style={black!50},
                grid=major, grid style={black!12},
                x tick label style={/pgf/number format/1000 sep={\,},
                        /pgf/number format/fixed},
                legend style={font=\scriptsize, draw=black!30, fill=white,
                        at={(0.98,0.03)}, anchor=south east},
            ]
            \addplot[azul, thick, mark=none] table[x=carga, y=Det] {\figDET/data/carga.dat};
            \addlegendentry{AST-Def-DETR}
            \addplot[ambar!90!black, thick, dashed, mark=none] table[x=carga, y=Frc] {\figDET/data/carga.dat};
            \addlegendentry{Faster R-CNN}
            \addplot[verde!60!black, very thick, mark=none] table[x=carga, y=Yol] {\figDET/data/carga.dat};
            \addlegendentry{YOLO26s}
            \addplot[rosa!80!black, thick, dotted, mark=none] table[x=carga, y=Rtd] {\figDET/data/carga.dat};
            \addlegendentry{RT-DETR-L}
            \draw[black!55, dashed] (axis cs:3000,0.2) -- (axis cs:3000,0.86)
            node[above, font=\scriptsize, text=black!65, align=center]
            {carga de los\\puntos de operación};
        \end{axis}
    \end{tikzpicture}
\end{figure}

Cada curva es un modelo: cuanto más alta, más eventos recupera con la misma
cantidad de propuestas por hora, y la línea punteada marca la carga de los cuatro
puntos de operación. La Tabla~\ref{tab:carga-igual} lee las curvas en seis
presupuestos.

\begin{table}[htbp]
    \centering
    \small
    \begin{threeparttable}
        \caption{Cobertura a carga igual}
        \label{tab:carga-igual}
        \begin{tabular}{@{}rcccc@{}}
            \toprule
            \thead{Propuestas\\por hora} & \thead{AST-Def-\\DETR} & \thead{Faster\\R-CNN} & \thead{YOLO26s} & \thead{RT-DETR-L} \\
            \midrule
            \input{\figDET/tablas/carga} \\
            \bottomrule
        \end{tabular}
        \begin{tablenotes}[flushleft]
            \footnotesize
            \item \textit{Nota.} Cada columna es la curva de un modelo leída a la misma abscisa.
            El techo es el \textit{recall} con el umbral en el piso del
            protocolo, esto es, con todo lo que el modelo es capaz de proponer;
            la última fila es la carga a la que lo alcanza. Un guion significa
            que el modelo no llega a emitir tantas propuestas por hora.
        \end{tablenotes}
    \end{threeparttable}
\end{table}

El resultado refuerza la lectura de la Tabla~\ref{tab:resultados-op} y añade algo
que el punto de operación no dejaba ver.

\rotulo{En el rango de carga que un revisor humano puede absorber, las cuatro
    configuraciones son indistinguibles.} A \rBudgetCommon{} propuestas por hora
(aproximadamente la carga de los cuatro puntos de operación, que promedian
\rMeanLoad{}) la cobertura va de \rCovkFrc{} a \rCovkYol{}. Cuatro
arquitecturas con \cParamsYol{}, \cParamsRtd{}, \cParamsFrc{} y \cParamsDet{}
millones de parámetros, con extractores preentrenados en tres corpus distintos y
con recetas de entrenamiento que no están igualadas, convergen a la misma
cobertura cuando se les da el mismo presupuesto de revisión. \emph{La elección de
    arquitectura pesa menos, en esta tarea y con este material, que la elección del
    punto de operación.}

\rotulo{Lo que separa a las curvas es dónde termina cada una.} El techo de
cobertura (lo que el modelo recupera emitiendo absolutamente todo lo que puede
proponer) va de \rCeilingYol{} (YOLO26s) a \rCeilingFrc{} (Faster R-CNN). Pero
el precio de ese techo es abismalmente distinto: Faster R-CNN necesita
\rCeilingLoadFrc{} propuestas por hora para alcanzarlo y YOLO26s llega al suyo
con \rCeilingLoadYol{}, un orden de magnitud menos. Ninguna de esas cifras es
operativa (revisar ochenta mil cajas por hora de audio no es asistencia),
pero la comparación sí dice algo: las curvas de los otros tres modelos solo
superan a YOLO26s en una región de carga que nadie va a usar.

\rotulo{El umbral de 0,5 de fábrica no habría comparado nada.} El umbral que la
validación elige para YOLO26s es \rThresholdYol{}, cinco veces menor. Operarlo a
0,5 lo habría situado en la parte conservadora de su curva y a Faster R-CNN
(cuyo umbral elegido es \rThresholdFrc{}) en la permisiva, y el orden entre
los dos se habría invertido sin que ninguna propiedad de las arquitecturas
cambiara. Es la razón por la que el punto de operación de cada modelo debe
elegirse barriendo su umbral, y no heredando uno.

\subsection{Lectura por clase}

El agregado está dominado por dos clases: \rWindowClasses{} aportan el
\rShareTwoWindow{} de las cajas de prueba y son, por la fragmentación que
documenta la sección~\ref{sec:ventaneo}, las más fáciles del conjunto, porque sus
cajas ocupan la ventana entera. Por eso la Tabla~\ref{tab:recall-clase} reporta
también el promedio por clase, que trata igual a una clase de 37 cajas que a una
de dos mil.

\begin{table}[htbp]
    \centering
    \scriptsize
    \begin{threeparttable}
        \caption{Recall por clase}
        \label{tab:recall-clase}
        \begin{tabular}{@{}L{1.9cm}rcccc@{}}
            \toprule
            \textbf{Clase} & \thead{Cajas\\test} & \thead{AST-Def-\\DETR} & \thead{Faster\\R-CNN} & \thead{YOLO26s} & \thead{RT-DETR-L} \\
            \midrule
            \input{\figDET/tablas/porclase} \\
            \bottomrule
        \end{tabular}
        \begin{tablenotes}[flushleft]
            \footnotesize
            \item \textit{Nota.} En negrita, el mejor de las cuatro configuraciones en cada
            clase. Un asterisco marca las dos clases de ventana, que no entran en
            el promedio por clase.
        \end{tablenotes}
    \end{threeparttable}
\end{table}

El promedio por clase separa las configuraciones mucho más que el agregado:
\rMacroYol{} para YOLO26s, \rMacroRtd{} para RT-DETR-L, \rMacroDet{} para la
propuesta y \rMacroFrc{} para Faster R-CNN. Es el mismo orden que la mAP, y con
una distancia mayor: el modelo seleccionado le saca seis puntos de promedio por
clase a la propuesta mientras que en el agregado le saca uno y medio. La ventaja
no viene de las clases fáciles, viene de las escasas.

Tres patrones se leen en la tabla.

\rotulo{Hay un suelo que ninguna configuración cruza.} \rHardClasses{} quedan por
debajo de 0,5 en las cuatro, y son dos de las \rNScarce{} clases con 60 cajas de
prueba o menos. Cuando las cuatro arquitecturas fallan por igual sobre una clase,
la causa razonable es el material y no el diseño: son las clases que el umbral de
\minPairCount{} anotaciones dejó entrar por poco y cuyo desempeño la
sección~\ref{sec:limitaciones} anticipó como no atribuible a la arquitectura.

\rotulo{La escasez no lo explica todo.} \code{lw/ta} (91 cajas) llega a 0,934 y
\code{pt/sqc} (138) a 0,993 con el modelo seleccionado, por encima de clases cinco
veces más numerosas. Ambas ocupan bandas de frecuencia estrechas y bien
separadas, lo que sugiere que la separabilidad geométrica de la
sección~\ref{sec:geometria} pesa más que el volumen de ejemplos. En la otra
dirección, \code{as/bc} (647 cajas) no pasa de 0,672 en ninguna configuración: es
una llamada que convive con la clase de ventana \code{as/hc} de la misma especie
y queda enterrada en el coro.

\rotulo{Ninguna configuración gana en todas partes.} El modelo seleccionado es el
mejor en 9 de las \rNDetectionClasses{} clases de detección, RT-DETR-L en 7, la
propuesta en 4 y Faster R-CNN en 3. La propuesta gana justamente donde el
contexto largo importa (\code{as/bc} y \code{lw/trino}, dos clases que se
emiten en secuencia), que es lo que cabría esperar de una cabeza de conjuntos
con atención sobre toda la ventana.

\subsection{Ejemplos cualitativos}

La Figura~\ref{fig:cualitativas} muestra seis ventanas de prueba con las
anotaciones del equipo y las predicciones del modelo seleccionado a su umbral.
Las seis se eligieron por criterio explícito y no por aspecto: dos escenas
densas, dos ventanas con todos sus eventos recuperados, una con un evento que el
modelo no encontró y una con un falso positivo de alta confianza, todas de
grabaciones distintas.

\begin{figure}[htbp]
    \centering
    \caption{Ejemplos cualitativos}
    \label{fig:cualitativas}
    \includegraphics[width=\textwidth]{\figDET/detecciones_cualitativas.png}
\end{figure}

Las dos escenas densas ilustran los dos extremos del comportamiento del modelo
sobre el mismo tipo de material. En la primera recupera la mayoría de una serie
de llamadas cortas y repetidas, con alguna caja de más; en la segunda, una serie
equivalente de la misma clase en otra grabación, no emite ni una sola caja por
encima del umbral. La diferencia entre ambas no está en la densidad sino en la
relación señal-ruido: la segunda tiene el coro de fondo encima de la banda de la
llamada. Es un modo de fallo que conviene caracterizar, porque determina qué
fracción del material queda sin asistir.

\subsection{Selección del modelo y qué establece la comparación}
\label{sec:seleccion-modelo}

\rotulo{El modelo que pasa a RE2.4 es \rWinner{}} (corrida \rWinnerRun{}). La
regla declarada lo decide en tres pasos, y conviene dejar escrito cada uno:

\begin{enumerate}[nosep]
    \item La métrica primaria (\textit{recall} en prueba con su intervalo)
          \textbf{no separa} a ninguna de las seis parejas posibles. No elige.
    \item La secundaria (mAP@\pMatchIou{}) ordena las cuatro y pone a
          \rBestSecondary{} primero con \rMapYol{} frente a \rMapRtd{} del
          segundo. Elige.
    \item El desempate por costo de inferencia no hizo falta, pero apunta en la
          misma dirección: el modelo seleccionado es también el más barato
          (sección~\ref{sec:costo-inferencia}).
\end{enumerate}

\rotulo{La hipótesis de diseño no se confirma.} La propuesta queda tercera de
cuatro en las dos métricas: \rRecallDet{} de \textit{recall} frente a
\rRecallYol{} del modelo seleccionado, y \rMapDet{} de mAP@\pMatchIou{} frente a
\rMapYol{}. En la primaria la diferencia cae dentro de los intervalos de
confianza; en la secundaria, la propuesta empata con Faster R-CNN y queda por
detrás de YOLO26s y de RT-DETR-L. No supera de forma medible a ninguna línea
base, que es lo que la hipótesis afirmaba.

RT-DETR-L, que se incluyó para esto (sección~\ref{sec:baselines}), indica de qué
componente viene el resultado. Predice un conjunto de cajas igual que la
propuesta, pero desde un extractor preentrenado en imágenes, y con
\cParamsRtd{}\,M de parámetros frente a \cParamsDet{}\,M queda a la par en la
primaria (\rRecallRtd{}) y por encima en la secundaria (\rMapRtd{} frente a
\rMapDet{}). El contraste apunta a que la cabeza de conjuntos no es lo que falla:
lo que no aporta una ventaja medible es el preentrenamiento en audio, en esta
tarea y con este presupuesto de entrenamiento. Como las recetas de
entrenamiento no están igualadas, es una indicación y no una atribución
(sección~\ref{sec:discusion-oe2}).

\subsection{Costo de inferencia}
\label{sec:costo-inferencia}

El costo se midió sobre las mismas \rTimingWindows{} ventanas de prueba y en la
misma GPU (\rGpuName{}), incluyendo el preprocesamiento de cada arquitectura, la
supresión de no máximos y el tope de \mMaxDet{} cajas por ventana. El modelo
seleccionado consume \rGpuMinYol{} minutos de GPU por hora de audio; la
propuesta, \rGpuMinDet{}; RT-DETR-L, \rGpuMinRtd{}; Faster R-CNN, \rGpuMinFrc{}.
Son \rCostRatio{} veces entre los extremos. En términos operativos, el modelo
seleccionado procesa el audio unas \rRealtimeYol{} veces más rápido que en tiempo
real y el más lento unas \rRealtimeFrc{} veces: las cuatro son viables para
monitoreo acústico pasivo, de modo que el costo de inferencia es un criterio de
desempate y no una restricción.

\subsection{Verificación de RE2.3}
\label{sec:verificacion-re23}

El indicador enumera siete magnitudes mínimas, exige separar los tres ejes y
justificar la selección del modelo. La Tabla~\ref{tab:verificacion-re23} las
localiza una por una.

\begin{table}[htbp]
    \centering
    \small
    \begin{threeparttable}
        \caption{Verificación de RE2.3}
        \label{tab:verificacion-re23}
        \begin{tabular}{@{}L{4.4cm}L{4.0cm}L{4.8cm}@{}}
            \toprule
            \textbf{Elemento del IOV} & \textbf{Dónde} & \textbf{Valor del modelo seleccionado} \\
            \midrule
            \textit{Recall} & Tabla~\ref{tab:resultados-op} & \rRecallYol{} \rRecallCiYol{} \\
            Precisión & Tabla~\ref{tab:resultados-op} & \rPrecisionYol{} \\
            F\textsubscript{1} & Tabla~\ref{tab:resultados-op} & \rFoneYol{} \\
            IoU medio & Tabla~\ref{tab:ejes}, Figura~\ref{fig:iou} & \rIouYol{} (mediano) \\
            AP a IoU 0,25 & Tabla~\ref{tab:resultados-op} & \rMapLooseYol{} \\
            AP a IoU 0,5 & Tabla~\ref{tab:resultados-op} & \rMapStrictYol{} \\
            \textit{Recall} por clase & Tabla~\ref{tab:recall-clase} & \rMacroYol{} de promedio por clase \\
            Matriz de confusión & Figura~\ref{fig:confusion} & \rConfusionOff{} fuera de la diagonal \\
            Separa detección, encuadre y clasificación & Sección~\ref{sec:ejes}, Tabla~\ref{tab:ejes} & los tres ejes, por separado \\
            Justifica la selección del modelo & Sección~\ref{sec:seleccion-modelo} & regla declarada, aplicada en tres pasos \\
            \bottomrule
        \end{tabular}
        \begin{tablenotes}[flushleft]
            \footnotesize
            \item \textit{Nota.} Los archivos de métricas son
            \file{comparacion\_\allowbreak modelos\_\allowbreak v3.txt}
            y su \file{.json} con las mismas cifras, en la carpeta RE2.3 del
            Anexo~\ref{anx:medios}. Toda magnitud de este
            capítulo procede de volver a aplicar el protocolo sobre los volcados
            crudos, no de una cifra transcrita.
        \end{tablenotes}
    \end{threeparttable}
\end{table}

Cada fila apunta a una tabla, una figura o una sección de este capítulo, y la
última columna da el valor del modelo seleccionado, de modo que el indicador
puede comprobarse sin salir del documento.

% ===================================================================
\section{Modelo final cargable y demostración funcional (RE2.4)}
\label{sec:re24}
% ===================================================================

\rotulo{Resultado.} El \textit{checkpoint} del modelo seleccionado, con su
vocabulario y su punto de operación, y la demostración de que una inferencia
sobre un WAV produce una tabla de selección de Raven con coordenadas
tiempo--frecuencia, especie, tipo de llamada y puntuación.

\rotulo{Cómo se alcanzó.} Es la Fase~6 de CRISP-DM (sección~\ref{sec:fase6}).
\file{detect.py} corta cada grabación en las mismas
ventanas de \clipLen{}\,s con salto de \clipHop{}\,s del caché, ejecuta el modelo
por lotes, aplica el posprocesamiento de la evaluación, funde entre ventanas
solapadas las cajas anidadas de la misma clase y escribe la tabla junto a cada
audio.

\rotulo{Cómo reproducirlo.} Basta el \textit{checkpoint} publicado en la versión
v0.2.0 y el código del repositorio: \file{src/detect.py} acepta un WAV o una
carpeta y escribe la tabla junto a cada audio. La demostración de la
sección~\ref{sec:demo-oe2} se obtuvo así sobre una grabación de prueba.

\rotulo{Medio de verificación.} El \textit{checkpoint} del modelo final, que lleva
dentro su vocabulario, y su punto de operación, publicados en la versión v0.2.0
del repositorio (\url{\rRepoUrl/releases/tag/v0.2.0}) como el archivo
\file{detector-2.0.0-model-yolo26s\_v3.zip}, con su suma de verificación SHA-256;
y la tabla de selección que la inferencia produce sobre una grabación de prueba
(sección~\ref{sec:demo-oe2}). El indicador nombra un archivo \file{.pth} y
\file{src/infer.py}, nombres del prototipo: en la implementación final el
\textit{checkpoint} es un \file{.pt} autocontenido y la inferencia la ejecuta
\file{src/detect.py}, con la misma función. La carpeta RE2.4 de los medios de
verificación (Anexo~\ref{anx:medios}) contiene la comprobación de que el
\textit{checkpoint} carga sin claves incompatibles, su suma SHA-256, su punto de
operación y la tabla de salida de la demostración.

\subsection{El modelo final}
\label{sec:modelo-final}

El modelo que pasa a producción es el ganador de la comparación: \rWinnerRun{},
la época \rFinalEpoch{} de su corrida, con \cParamsYol{}\,M de parámetros en un
\textit{checkpoint} de \rFinalMb{}\,MB. Su punto de operación, elegido en
validación con precisión $\geq$ \pMinPrecision{}, es un umbral de confianza de
\rFinalThreshold{}, al que en prueba recupera \rRecallYol{} de las cajas
anotadas con precisión \rPrecisionYol{}. El umbral viaja junto al
\textit{checkpoint} en \file{operating\_point.json} y es el que \file{detect.py}
usa por defecto, de modo que el usuario no tiene que conocerlo.

\rotulo{Carga sin claves incompatibles.} \code{load\_checkpoint} reconstruye la
arquitectura a partir de las claves \code{architecture} y \code{hparams} del
archivo, con el número de clases que dicta \code{labels} (\rFinalClasses{}), y
carga el \code{state\_dict} en modo estricto: los \rFinalTensors{} tensores del
archivo coinciden uno a uno con los del modelo reconstruido, sin faltantes ni
sobrantes.

\subsection{Formato de salida}
\label{sec:formato-salida}

La salida es una \emph{tabla de selección} de Raven Pro: texto delimitado por
tabulaciones con las \rDemoColumns{} columnas de la Tabla~\ref{tab:columnas-salida}, que
Raven abre sobre la grabación y dibuja como cajas sobre el espectrograma. Es el
mismo formato en el que el equipo entrega sus anotaciones
(sección~\ref{sec:material}), con dos diferencias: la columna \code{Score}, que
Raven ignora y permite ordenar o filtrar las detecciones, y que \code{Species} y
\code{Call type} vienen siempre en mayúscula y del vocabulario cerrado. Las
coordenadas de frecuencia se obtienen invirtiendo la escala mel con la que se
normalizaron las cajas (ecuación~\ref{eq:mel-inv}).

Que la salida tenga esta forma no es un detalle de implementación: es la tercera
forma de salida de la sección~\ref{sec:soluciones}, la única que coincide con la
anotación del especialista, y llega en el archivo que el analista ya usa.

\begin{table}[htbp]
    \centering
    \small
    \caption{Columnas de la tabla de salida}
    \label{tab:columnas-salida}
    \begin{tabular}{@{}lL{9.4cm}@{}}
        \toprule
        \textbf{Columna} & \textbf{Contenido} \\
        \midrule
        \code{Selection} & Correlativo desde 1, como exige Raven \\
        \code{View}, \code{Channel} & \code{Spectrogram 1}, \code{1}: la vista y el canal en los que Raven dibuja \\
        \code{Begin Time (s)}, \code{End Time (s)} & Lados temporales de la caja, en segundos desde el inicio del archivo \\
        \code{Low Freq (Hz)}, \code{High Freq (Hz)} & Lados de frecuencia, en hercios \\
        \code{Species} & Código de especie en mayúscula (\code{LW}) \\
        \code{Call type} & Código del tipo de llamada en mayúscula (\code{CS}) \\
        \code{Score} & Confianza del modelo en $[0, 1]$; solo en tablas del modelo \\
        \bottomrule
    \end{tabular}
\end{table}

Las columnas de coordenadas y de etiquetas son las mismas que escribe el analista;
\code{Score} es la única que añade el modelo. La
Figura~\ref{fig:flujo-inferencia} resume cómo se llega a ellas desde la grabación.

\begin{figure}[htbp]
    \centering
    \caption{Flujo de inferencia}
    \label{fig:flujo-inferencia}
    \resizebox{\textwidth}{!}{%
        \begin{tikzpicture}[node distance=4mm and 7mm,
                caja/.append style={font=\scriptsize, inner sep=3pt},
                dato/.append style={font=\scriptsize\ttfamily}]
            \node[dato] (wav) {grabación.wav};
            \node[etapa, right=of wav] (win) {ventanas \clipLen{}\,s / \clipHop{}\,s\\mel de cada una};
            \node[etapa, right=of win] (model) {modelo por lotes\\score $\geq$ umbral};
            \node[etapa, right=of model] (post) {NMS del modelo\\tope \mMaxDet{} por ventana};
            \node[etapa, below=of post] (merge)
            {fusión entre ventanas\\anidadas de la misma clase, IoU \mNmsIou{}};
            \node[etapa, left=of merge] (hz) {tiempo absoluto\\mel $\to$ Hz};
            \node[dato, left=of hz] (txt) {grabación.detections.txt};
            \draw[flecha] (wav) -- (win);
            \draw[flecha] (win) -- (model);
            \draw[flecha] (model) -- (post);
            \draw[flecha] (post) -- (merge);
            \draw[flecha] (merge) -- (hz);
            \draw[flecha] (hz) -- (txt);
            \node[nota, below=1mm of txt] {Raven Pro $|$ hoja de cálculo};
        \end{tikzpicture}}
\end{figure}

El paso que no existe en la evaluación es la fusión entre ventanas: como las
ventanas se solapan, una misma llamada puede detectarse dos veces, y la fusión
deja una sola caja antes de pasar a tiempo absoluto y a hercios.

\subsection{Demostración}
\label{sec:demo-oe2}

El modelo final se aplicó sobre una grabación de la partición de prueba, que no
participó del entrenamiento ni de la selección del \textit{checkpoint}:
\rDemoFile{}, de \rDemoSeconds{} segundos, con \rDemoAnnotations{} anotaciones
del equipo repartidas entre \rDemoClasses{}. La inferencia corrió en CPU, que es
la condición más desfavorable y la que tendría un analista sin GPU, y produjo una
tabla con las \rDemoColumns{} columnas de la Tabla~\ref{tab:columnas-salida} y una fila
por detección. La Tabla~\ref{tab:salida} recoge sus primeras filas.

\begin{table}[htbp]
    \centering
    \small
    \caption{Primeras filas de la tabla de salida}
    \label{tab:salida}
    \begin{tabular}{@{}rrrrrllr@{}}
        \toprule
        \thead{Sel.} & \thead{Begin\\Time (s)} & \thead{End\\Time (s)} & \thead{Low\\Freq (Hz)} & \thead{High\\Freq (Hz)} & \thead{Species} & \thead{Call\\type} & \thead{Score} \\
        \midrule
        \input{\figDET/tablas/salida} \\
        \bottomrule
    \end{tabular}
\end{table}

Sobre esa grabación el modelo emite \rDemoDetections{} detecciones al umbral
\rFinalThreshold{}. A IoU $\geq$ \pMatchIou{} en coordenadas de la grabación
entera, \rDemoFound{} de las \rDemoAnnotations{} anotaciones del equipo tienen
una detección encima, \rDemoHits{} de las \rDemoDetections{} detecciones caen
sobre una anotación y \rDemoSameLabel{} de ellas llevan además la misma etiqueta.
La Figura~\ref{fig:demo} superpone las dos capas, tal como se ven en Raven,
sobre el tramo de \rDemoWindow{} segundos que concentra más
anotaciones. \textbf{Es un ejemplo, no una medición}: una grabación no es una
muestra, y las métricas son las de la sección~\ref{sec:re23}.

\begin{figure}[htbp]
    \centering
    \caption{Demostración sobre una grabación}
    \label{fig:demo}
    \includegraphics[width=\textwidth]{\figDET/demo_grabacion.png}
\end{figure}

El archivo se abre en Raven sobre la grabación como una tabla de selección más,
sin conversión ni herramienta añadida. La Figura~\ref{fig:raven-compat} lo
muestra sobre otra grabación de prueba, \file{20240117\_162715.wav} de
\textit{Leontocebus weddelli}: Raven Lite carga \file{20240117\_162715.detections.txt}
como \emph{Selection Table}, dibuja las seis detecciones como cajas en sus
coordenadas de tiempo y frecuencia, las numera sobre el espectrograma y arrastra
la columna \code{Species} a la tabla. A las columnas que produjo el modelo añade
las que calcula por su cuenta (\code{Delta Time}, \code{Delta Freq},
\code{Avg Power Density}), igual que con una tabla anotada a mano. \code{Call
    type} y \code{Score} viajan en el archivo pero esta versión del programa no
las muestra, de modo que el filtrado por confianza se hace fuera de Raven.

\begin{figure}[htbp]
    \centering
    \caption{La tabla generada abierta en Raven}
    \label{fig:raven-compat}
    \includegraphics[width=\textwidth]{\figDET/raven_compatibility.png}
\end{figure}

A partir de ahí el trabajo es el habitual: el analista oye cada caja, borra las
que no correspondan, corrige las que estén mal encuadradas y guarda la tabla ya
revisada. Ese es el circuito cuyo tiempo mide OE3.

\subsection{Verificación de RE2.4}

El indicador exige dos cosas, y las dos quedan demostradas arriba:

\begin{enumerate}[nosep]
    \item \textbf{El \textit{checkpoint} se carga sin claves incompatibles.} La
          carga en modo estricto de la sección~\ref{sec:modelo-final} no deja
          tensores faltantes ni sobrantes.
    \item \textbf{La inferencia sobre un WAV produce un archivo tabulado con las
              columnas de Raven.} La Tabla~\ref{tab:salida} lo muestra:
          \rDemoDetections{} filas con coordenadas tiempo--frecuencia, especie,
          tipo de llamada y puntuación, en las \rDemoColumns{} columnas de la
          Tabla~\ref{tab:columnas-salida}. Que esas columnas son las de Raven y no
          una imitación lo comprueba la Figura~\ref{fig:raven-compat}: el archivo
          se abre en el propio programa, sobre su grabación, y las detecciones
          quedan dibujadas sobre el espectrograma.
\end{enumerate}

% ===================================================================
\section{Código y modelo disponibles en un repositorio (RE2.5)}
\label{sec:re25}
% ===================================================================

\rotulo{Resultado.} El repositorio público \url{\rRepoUrl}, que contiene el
código de curación, entrenamiento, evaluación e inferencia y las instrucciones de
ejecución, y que publica el \textit{checkpoint} del modelo final en su versión
v0.2.0.

\rotulo{Cómo se alcanzó.} Con el procedimiento de respaldo de la
sección~\ref{sec:procedimientos}, versionando el proyecto desde el inicio y no al
entregarlo: \rNCommits{} \textit{commits} y una versión declarada en
\file{pyproject.toml} (\rVersion{}), con la trazabilidad que describe la
sección~\ref{sec:reproducibilidad}.

\rotulo{Cómo reproducirlo.} Por dos caminos. Desde el código, clonando el
repositorio en el \textit{commit} \rCommit{}, instalando el entorno de
\file{uv.lock} y ejecutando las etapas de la sección~\ref{sec:re22}. Sin instalar
nada, con el \file{README} del repositorio, que explica cómo descargar de su
página de versiones el programa para Windows y el modelo, y cómo procesar con
ellos una grabación.

\rotulo{Medio de verificación.} La URL del repositorio y la de su versión
v0.2.0 (\url{\rRepoUrl/releases/tag/v0.2.0}), que publica el modelo final; esta
sección es el índice que lleva de cada verbo del indicador al archivo que lo
implementa, y la carpeta RE2.5 de los medios de verificación
(Anexo~\ref{anx:medios}) lo reproduce.

\subsection{Estructura}

El repositorio separa el código del material que no cabe en él. La
Figura~\ref{fig:arbol} muestra el primer nivel: \file{src/} contiene el
\textit{pipeline} entero y \file{research/}, este documento con sus figuras. Las
dos carpetas sin versionar son las que guardan datos del equipo de investigación
y artefactos de decenas de gigabytes, que se reconstruyen con el código que sí
está publicado.

\begin{figure}[htbp]
    \centering
    \caption{Árbol del repositorio}
    \label{fig:arbol}
    \begin{minipage}{0.92\textwidth}
\begin{lstlisting}
tesis-primate/
  README.md         que es el proyecto y como se ejecuta
  pyproject.toml    dependencias, extras y grupos; uv.lock las fija
  evaluation.sh     volcado + comparacion de todas las corridas
  Makefile          compila la tesis
  src/              el pipeline entero (figura 5.9)
  notebooks/        auditoria, curacion asistida, preentrenamiento con aves
  research/         la tesis: main.tex, capitulos/<nn>/{tex, figures, notebook}
  docs/             requisitos del sistema y notas técnicas
  data/   (no versionado)  raw/, cleaned/, processed/, yolo/
  runs/   (no versionado)  corridas: checkpoints, volcados, comparacion
\end{lstlisting}
    \end{minipage}
\end{figure}

El indicador de RE2.5 enumera lo que el repositorio debe permitir hacer. La
Tabla~\ref{tab:donde} lleva cada uno de esos verbos al archivo que lo implementa
y al resultado esperado que lo documenta.

\begin{table}[htbp]
    \centering
    \small
    \caption{Ubicación de cada elemento del indicador}
    \label{tab:donde}
    \begin{tabular}{@{}L{3.0cm}L{7.0cm}L{3.4cm}@{}}
        \toprule
        \textbf{Qué} & \textbf{Dónde} & \textbf{Documentado en} \\
        \midrule
        Curación & \file{src/prepare\_annotations.py}, \file{src/data/annotations.py}, \file{src/data/species.py} & RE1.1 \\
        Conjunto y partición & \file{src/prepare\_data.py}, \file{src/data/manifest.py}, \file{src/data/cache.py} & RE1.2, RE1.3 \\
        Arquitecturas & \file{src/models/} & RE2.1 \\
        Entrenamiento & \file{src/train.py}, \file{src/training/} & RE2.2 \\
        Evaluación & \file{src/dump\_predictions.py}, \file{src/compare\_models.py}, \file{src/evaluation/}, \file{evaluation.sh} & RE2.1, RE2.3 \\
        Inferencia & \file{src/detect.py}, \file{src/inference/} & RE2.4 \\
        Instrucciones de ejecución & \file{README.md}, \file{pyproject.toml}, Listado~\ref{lst:pipeline} & RE2.2 \\
        Modelo final & El \textit{checkpoint} autocontenido de \rWinnerRun{} con su \file{operating\_point.json}, en la versión v0.2.0 & Sección~\ref{sec:modelo-final} \\
        \bottomrule
    \end{tabular}
\end{table}

Cada fila apunta a una ruta del repositorio público, salvo la última, que es un
archivo de su versión v0.2.0.

\subsection{Qué contiene y qué no}

\rotulo{El código, completo.} Todo el \textit{pipeline} tal como lo describe la
sección~\ref{sec:re22}, desde la lectura de las tablas de Raven hasta la escritura
de la tabla de salida. No depende de nada que no esté declarado en
\file{pyproject.toml} y fijado en \file{uv.lock}.

\rotulo{El modelo final, no las corridas.} El artefacto que se publica es el
\textit{checkpoint} de \rWinnerRun{} (autocontenido, con su arquitectura, sus
hiperparámetros, su vocabulario y el \file{meta.json} del caché con el que se
entrenó (Tabla~\ref{tab:checkpoint})) junto con su punto de operación. Las
corridas completas, con sus \file{last.pt} y sus volcados, suman decenas de
gigabytes y no aportan nada que el \textit{checkpoint} elegido y las cifras de
este capítulo no digan ya.

\rotulo{Las grabaciones y las anotaciones primarias, no.} Pertenecen al equipo de
investigación. El repositorio no las contiene ni las referencia por ubicación, y
los pesos preentrenados de terceros (el AST de AudioSet, los de COCO) se
resuelven desde su origen y conservan la licencia de su autor.

\rotulo{Fuera de alcance.} La distribución del código y del modelo como producto
(empaquetado para usuarios finales, versionado para terceros, licenciamiento y
mantenimiento) no forma parte de esta tesis. El compromiso de este resultado,
tal como lo enuncia su indicador, es que la URL permita localizar el código y los
artefactos del modelo final, y eso es lo que la Tabla~\ref{tab:donde} verifica.

\subsection{Verificación de RE2.5}

El indicador pide que la URL permita localizar cuatro cosas, y las cuatro quedan
localizadas en la Tabla~\ref{tab:donde}: el código de curación, el de
entrenamiento, el de evaluación y el de inferencia, cada uno con el resultado
esperado que lo documenta; las instrucciones de ejecución, en el
Listado~\ref{lst:pipeline}; y los artefactos del modelo final, que son el
\textit{checkpoint} de \rWinnerRun{} y su punto de operación.

% ===================================================================
\section{Discusión}
\label{sec:discusion-oe2}
% ===================================================================

\rotulo{Qué se obtuvo.} Se implementó y se evaluó un detector que recibe un
espectrograma y devuelve cajas tiempo--frecuencia con especie, tipo de llamada y
confianza, y se lo comparó con tres líneas base bajo un protocolo declarado antes
de correr el experimento. Sobre las \nWindowsTest{} ventanas de prueba, el modelo
seleccionado recupera \rRecallYol{} de las \nBoxesTest{} cajas anotadas
\rRecallCiYol{} con precisión \rPrecisionYol{}, al umbral que la validación
eligió; como la métrica primaria no separa a las cuatro configuraciones, lo
eligió la secundaria (sección~\ref{sec:seleccion-modelo}). De los \rDemoDetections{} resultados que produce sobre una
grabación de \rDemoSeconds{} segundos, el sistema entrega una tabla que el
analista abre en su propia herramienta.

\rotulo{Qué significan.} Tres lecturas, en orden de importancia para el problema.

La primera es que \textbf{el error está en encontrar y encuadrar, no en
    nombrar}. Una vez localizado el evento, la clase casi nunca falla
(sección~\ref{sec:ejes}): la parte del problema que parecía más difícil, con
\nClasses{} clases desbalanceadas, resultó ser la que el modelo resuelve. Para el uso asistido esto es una buena noticia con
una consecuencia concreta: el revisor no tendrá que corregir etiquetas, tendrá
que confirmar o descartar cajas, que es una operación mucho más rápida. Y para el
diseño de un sistema futuro dice dónde invertir: en la detección y el encuadre, no
en la cabeza de clasificación.

La segunda es que \textbf{la elección del punto de operación pesa más que la
    elección de la arquitectura}. A \rBudgetCommon{} propuestas por hora de audio
las cuatro configuraciones cubren entre \rCovkFrc{} y \rCovkYol{}, pero mover el
umbral de un solo modelo lo lleva de 0,52 a 0,91 dentro del mismo rango de carga.
Esto reordena la prioridad del trabajo que queda: una vez que se tiene un
detector razonable, el margen está en calibrarlo contra el presupuesto de revisión
real, que es justamente lo que mide OE3.

La tercera es que \textbf{el preentrenamiento en audio no aportó una ventaja
    medible}. Se esperaba que la propuesta superara a los detectores
preentrenados en COCO, y quedó tercera de cuatro
(sección~\ref{sec:seleccion-modelo}). Para el equipo es la lectura favorable: el
modelo seleccionado es también el más pequeño (\cParamsYol{}\,M de parámetros
frente a \cParamsDet{}\,M), el que menos entrenamiento necesitó (\cHoursYol{}\,h
de GPU frente a \cHoursDet{}\,h) y el más rápido en inferencia. OE2 pide un
detector comparado con una línea base, no que gane la propuesta, y la
herramienta usa el mejor de los cuatro. Para la investigación es
un resultado negativo, y vale como tal porque la regla de decisión se fijó antes
de entrenar. La propuesta gana algunas comparaciones parciales (cuatro clases, y
las dos clases que se emiten en secuencia), pero elegirlas después de ver los
resultados equivaldría a cambiar esa regla.

\rotulo{Consistencia con la literatura.} Los resultados concuerdan con la
línea del Capítulo~\ref{cap:estado} en tres puntos y discrepan en uno.

Concuerdan en que \textbf{la detección bidimensional sobre el espectrograma es
    una formulación viable en bioacústica}: \citet{escobar-amado_computer_2024}
reportan con Faster R-CNN sobre llamadas de foca, y \citet{airale_nbm_2025} y
\citet{gonzalez_yolo-based_2025} con detectores de la familia YOLO, cifras del
mismo orden que las de la Tabla~\ref{tab:resultados-op} sobre repertorios de
tamaño comparable. Concuerdan también en que \textbf{un umbral de IoU bajo es lo
    razonable cuando las fronteras del evento son borrosas} \citep{zhu_sound_2025}:
la Figura~\ref{fig:iou} muestra que incluso los pares que el protocolo cuenta
como aciertos tienen un IoU mediano de alrededor de 0,78, lejos de lo que se
exige en detección sobre imágenes naturales. Y concuerdan en que \textbf{la
    salida debe tener la forma que el analista ya usa}: la diferencia entre entregar
una probabilidad por segmento \citep{kahl_birdnet_2021} y entregar una caja con
especie, tipo de llamada y confianza es lo que permite abrir el resultado en
Raven sin ningún paso intermedio.

Discrepan con \citet{zhu_sound_2025}, el precedente directo de esta arquitectura,
que con un AST \emph{autosupervisado} supera a Faster R-CNN en un dominio de una
sola clase de sonido. Aquí el extractor es supervisado y las clases son
\nClasses{}, desbalanceadas: este trabajo no refuta aquel resultado, muestra que
no se traslada sin más a un repertorio de primates.

\rotulo{Generalización.} Qué de esto sobrevive fuera de este corpus.

Se generaliza, con bastante confianza, \textbf{la separación de los tres ejes
    como método de evaluación}. Es independiente del taxón y del repertorio, y su
valor (saber si un agregado bajo se debe a detección o a clasificación) se
vuelve mayor cuanto más clases tiene el problema. Se generaliza también
\textbf{la lectura de cobertura frente a carga}: cualquier sistema de
pre-anotación asistida compite contra un presupuesto de revisión, y el umbral
compartido entre arquitecturas es un artefacto de calibración en todas partes,
no solo aquí.

Se generaliza con reservas \textbf{el hallazgo de que la clasificación es fácil}.
Depende de que las clases sean acústicamente separables, y en este corpus lo son
en buena medida porque la especie está determinada por la carpeta de origen
(sección~\ref{sec:regla-especie}): un modelo que vea una grabación de una especie
no anotada por el equipo no tiene por qué comportarse igual. Un corpus con varias
especies simpátricas anotadas en la misma grabación pondría a prueba ese
resultado.

No se generaliza \textbf{el orden entre arquitecturas}. Es el resultado de este
material, con este presupuesto de entrenamiento y con recetas no igualadas; con
diez veces más datos anotados, o con un programa de entrenamiento más largo para
la propuesta, el orden podría ser otro. Tampoco se generalizan los umbrales
elegidos, que son propiedades de cada \textit{checkpoint} y no de la arquitectura.

\rotulo{Limitaciones.} Cinco, ordenadas por cuánto comprometen las conclusiones.

\begin{enumerate}
    \item \textbf{No hay techo de evaluación.} No se midió el acuerdo entre
          anotadores sobre una muestra del material, de modo que no se sabe cuánto
          del error que la Tabla~\ref{tab:resultados-op} atribuye al modelo es en
          realidad desacuerdo entre humanos. Es la limitación más seria de todo el
          capítulo: sin ella no se puede decir si \rRecallYol{} está cerca o lejos
          de lo alcanzable. OE3 la aborda parcialmente al hacer que un experto
          revise una muestra de predicciones.
    \item \textbf{La comparación no aísla componentes.} Las configuraciones
          usan el optimizador y las aumentaciones de su marco de trabajo, de modo
          que la comparación mide sistemas completos, que es lo que se declaró, y
          no permite atribuir la diferencia a un componente. Las dos variantes de
          la propuesta que lo harían (extractor congelado, y Deformable DETR con
          su ResNet-50 de COCO entrenado con la misma receta) ya están en el
          código y quedan como trabajo futuro.
    \item \textbf{Las clases escasas no se evalúan de forma fiable.}
          \rNScarce{} de las \nClasses{} clases tienen 60 cajas de prueba o menos,
          y sobre esas cifras un intervalo de confianza por clase sería tan ancho
          que no distinguiría nada. Las cifras por clase de la
          Tabla~\ref{tab:recall-clase} son indicativas en la mitad inferior de la
          tabla, no concluyentes.
    \item \textbf{Las dos clases de ventana inflan los agregados.} En
          \rWindowClasses{} detectar equivale a clasificar la ventana, y aportan
          el \rShareTwoWindow{} de las cajas de prueba: todo agregado que las
          incluya, como el \textit{recall} global, sale más alto que el de las
          clases de evento.
    \item \textbf{La prueba se midió una vez, por diseño, pero con una sola
              semilla.} El protocolo protege contra el sobreajuste a la partición de
          prueba, no contra la varianza de entrenamiento: no se sabe cuánto se
          movería cada cifra al reentrenar la misma configuración con otra semilla.
          Dado que las cuatro caben en \rRecallSpread{}, esa varianza podría ser del
          mismo orden que las diferencias que se reportan.
\end{enumerate}
